% Options for packages loaded elsewhere
\PassOptionsToPackage{unicode}{hyperref}
\PassOptionsToPackage{hyphens}{url}
\documentclass[
]{article}
\usepackage{xcolor}
\usepackage{amsmath,amssymb}
\setcounter{secnumdepth}{-\maxdimen} % remove section numbering
\usepackage{iftex}
\ifPDFTeX
  \usepackage[T1]{fontenc}
  \usepackage[utf8]{inputenc}
  \usepackage{textcomp} % provide euro and other symbols
\else % if luatex or xetex
  \usepackage{unicode-math} % this also loads fontspec
  \defaultfontfeatures{Scale=MatchLowercase}
  \defaultfontfeatures[\rmfamily]{Ligatures=TeX,Scale=1}
\fi
\usepackage{lmodern}
\ifPDFTeX\else
  % xetex/luatex font selection
\fi
% Use upquote if available, for straight quotes in verbatim environments
\IfFileExists{upquote.sty}{\usepackage{upquote}}{}
\IfFileExists{microtype.sty}{% use microtype if available
  \usepackage[]{microtype}
  \UseMicrotypeSet[protrusion]{basicmath} % disable protrusion for tt fonts
}{}
\makeatletter
\@ifundefined{KOMAClassName}{% if non-KOMA class
  \IfFileExists{parskip.sty}{%
    \usepackage{parskip}
  }{% else
    \setlength{\parindent}{0pt}
    \setlength{\parskip}{6pt plus 2pt minus 1pt}}
}{% if KOMA class
  \KOMAoptions{parskip=half}}
\makeatother
\usepackage{color}
\usepackage{fancyvrb}
\newcommand{\VerbBar}{|}
\newcommand{\VERB}{\Verb[commandchars=\\\{\}]}
\DefineVerbatimEnvironment{Highlighting}{Verbatim}{commandchars=\\\{\}}
% Add ',fontsize=\small' for more characters per line
\newenvironment{Shaded}{}{}
\newcommand{\AlertTok}[1]{\textcolor[rgb]{1.00,0.00,0.00}{\textbf{#1}}}
\newcommand{\AnnotationTok}[1]{\textcolor[rgb]{0.38,0.63,0.69}{\textbf{\textit{#1}}}}
\newcommand{\AttributeTok}[1]{\textcolor[rgb]{0.49,0.56,0.16}{#1}}
\newcommand{\BaseNTok}[1]{\textcolor[rgb]{0.25,0.63,0.44}{#1}}
\newcommand{\BuiltInTok}[1]{\textcolor[rgb]{0.00,0.50,0.00}{#1}}
\newcommand{\CharTok}[1]{\textcolor[rgb]{0.25,0.44,0.63}{#1}}
\newcommand{\CommentTok}[1]{\textcolor[rgb]{0.38,0.63,0.69}{\textit{#1}}}
\newcommand{\CommentVarTok}[1]{\textcolor[rgb]{0.38,0.63,0.69}{\textbf{\textit{#1}}}}
\newcommand{\ConstantTok}[1]{\textcolor[rgb]{0.53,0.00,0.00}{#1}}
\newcommand{\ControlFlowTok}[1]{\textcolor[rgb]{0.00,0.44,0.13}{\textbf{#1}}}
\newcommand{\DataTypeTok}[1]{\textcolor[rgb]{0.56,0.13,0.00}{#1}}
\newcommand{\DecValTok}[1]{\textcolor[rgb]{0.25,0.63,0.44}{#1}}
\newcommand{\DocumentationTok}[1]{\textcolor[rgb]{0.73,0.13,0.13}{\textit{#1}}}
\newcommand{\ErrorTok}[1]{\textcolor[rgb]{1.00,0.00,0.00}{\textbf{#1}}}
\newcommand{\ExtensionTok}[1]{#1}
\newcommand{\FloatTok}[1]{\textcolor[rgb]{0.25,0.63,0.44}{#1}}
\newcommand{\FunctionTok}[1]{\textcolor[rgb]{0.02,0.16,0.49}{#1}}
\newcommand{\ImportTok}[1]{\textcolor[rgb]{0.00,0.50,0.00}{\textbf{#1}}}
\newcommand{\InformationTok}[1]{\textcolor[rgb]{0.38,0.63,0.69}{\textbf{\textit{#1}}}}
\newcommand{\KeywordTok}[1]{\textcolor[rgb]{0.00,0.44,0.13}{\textbf{#1}}}
\newcommand{\NormalTok}[1]{#1}
\newcommand{\OperatorTok}[1]{\textcolor[rgb]{0.40,0.40,0.40}{#1}}
\newcommand{\OtherTok}[1]{\textcolor[rgb]{0.00,0.44,0.13}{#1}}
\newcommand{\PreprocessorTok}[1]{\textcolor[rgb]{0.74,0.48,0.00}{#1}}
\newcommand{\RegionMarkerTok}[1]{#1}
\newcommand{\SpecialCharTok}[1]{\textcolor[rgb]{0.25,0.44,0.63}{#1}}
\newcommand{\SpecialStringTok}[1]{\textcolor[rgb]{0.73,0.40,0.53}{#1}}
\newcommand{\StringTok}[1]{\textcolor[rgb]{0.25,0.44,0.63}{#1}}
\newcommand{\VariableTok}[1]{\textcolor[rgb]{0.10,0.09,0.49}{#1}}
\newcommand{\VerbatimStringTok}[1]{\textcolor[rgb]{0.25,0.44,0.63}{#1}}
\newcommand{\WarningTok}[1]{\textcolor[rgb]{0.38,0.63,0.69}{\textbf{\textit{#1}}}}
\usepackage{longtable,booktabs,array}
\newcounter{none} % for unnumbered tables
\usepackage{calc} % for calculating minipage widths
% Correct order of tables after \paragraph or \subparagraph
\usepackage{etoolbox}
\makeatletter
\patchcmd\longtable{\par}{\if@noskipsec\mbox{}\fi\par}{}{}
\makeatother
% Allow footnotes in longtable head/foot
\IfFileExists{footnotehyper.sty}{\usepackage{footnotehyper}}{\usepackage{footnote}}
\makesavenoteenv{longtable}
\setlength{\emergencystretch}{3em} % prevent overfull lines
\providecommand{\tightlist}{%
  \setlength{\itemsep}{0pt}\setlength{\parskip}{0pt}}
\usepackage{bookmark}
\IfFileExists{xurl.sty}{\usepackage{xurl}}{} % add URL line breaks if available
\urlstyle{same}
\hypersetup{
  hidelinks,
  pdfcreator={LaTeX via pandoc}}

\author{}
\date{}

\begin{document}

\section{Lowering Entry Barriers to Developing Custom Simulators of
Distributed Applications and Platforms with
SimGrid*}\label{lowering-entry-barriers-to-developing-custom-simulators-of-distributed-applications-and-platforms-with-simgrid}

Henri Casanova \(^{1}\) , Arnaud Giersch \(^{2}\) , Arnaud Legrand
\(^{3}\) , Martin Quinson \(^{4}\) , Frédéric Suter \(^{5}\)

\(^{1}\) Information and Computer Sciences, University of Hawai'i at
Manoa, Honolulu, HI, USA \(^{2}\) Université Marie et Louis Pasteur,
CNRS, Institut FEMTO-ST, F-90000 Belfort, France \(^{3}\) Université
Grenoble Alpes, CNRS, Inria, Grenoble INP, LIG, Grenoble, France
\(^{4}\) Rennes University, Inria, IRISA, Rennes, France \(^{5}\) Oak
Ridge National Laboratory, Oak Ridge, TN, USA

\section{Abstract}\label{abstract}

Researchers in parallel and distributed computing (PDC) often resort to
simulation because experiments conducted using a simulator can be for
arbitrary experimental scenarios, are less resource-, labor-, and
time-consuming than their real-world counterparts, and are perfectly
repeatable and observable. Many frameworks have been developed to ease
the development of PDC simulators, and these frameworks provide
different levels of accuracy, scalability, versatility, extensibility,
and usability. The SimGrid framework has been used by many PDC
researchers to produce a wide range of simulators for over two decades.
Its popularity is due to a large emphasis placed on accuracy,
scalability, and versatility, and is in spite of shortcomings in terms
of extensibility and usability. Although SimGrid provides sensible
simulation models for the common case, it was difficult for users to
extend these models to meet domain-specific needs. Furthermore, SimGrid
only provided relatively low-level simulation abstractions, making the
implementation of a simulator of a complex system a labor-intensive
undertaking. In this work we describe developments in the last decade
that have contributed to vastly improving extensibility and usability,
thus lowering or removing entry barriers for users to develop custom
SimGrid simulators.

\section{1 Introduction}\label{introduction}

Many parallel and distributed computing (PDC) research results are
obtained, at least in part, based on experiments conducted in
simulation. Reasons for PDC researchers to use simulation are: the
ability to explore arbitrary experimental scenarios; the fact that
simulation experiments can be less labor-, time- and resource-intensive
than their real-world counterparts; and the fact that simulation
experiments can be precisely controlled and instrumented, making them
perfectly observable and repeatable. Many discrete-event simulation
frameworks have been developed to provide abstractions and models for
compute, communication, and I/O resources and their usage, thus easing
the implementation of PDC simulators. The main concerns for these
framework are: (i) accuracy -- do their simulation models make it
possible for the simulated behavior to match that of the real-world
system being simulated? (ii) scalability -- do they allow the simulation
of large and long-running scenarios with low computational complexity
and memory footprint? (iii) versatility -- can they be used to develop
simulators for a diverse range of PDC scenarios and domains? (iv)
extensibility -- can their simulation models be configured, extended, or
even replaced so that the simulator can serve specific purposes? and (v)
usability -- do they make it possible to implement simulators with low
software engineering and development effort?

PDC simulation frameworks have been developed for decades, placing
different levels of emphasis on and using different approaches for
achieving compromises between these often conflicting concerns. For
instance, a possible approach for achieving higher accuracy is to
implement the simulation at a higher level of detail, but doing so
typically increases computational complexity and memory footprint, thus
reducing scalability. A way to increase scalability is to target a
specific PDC domain, which makes it possible to eschew simulating
features of the real-world system that are not relevant to that specific
domain. For instance, when simulating an IoT system in which only small
messages are exchanged, one can forgo the simulation of network
contention effects and still produce meaningful results. But doing so
reduces versatility by design.

The SimGrid framework {[}1{]} has demonstrated that it is

\textit{[Image omitted: images/7a7f4c4f936af5b3024e2c8854e511ddac282cf61c967189a0ec549e3f658998.jpg]}

bar\_stacked

{\def\LTcaptype{none} % do not increment counter
\begin{longtable}[]{@{}
  >{\raggedright\arraybackslash}p{(\linewidth - 12\tabcolsep) * \real{0.1429}}
  >{\raggedright\arraybackslash}p{(\linewidth - 12\tabcolsep) * \real{0.1429}}
  >{\raggedright\arraybackslash}p{(\linewidth - 12\tabcolsep) * \real{0.1429}}
  >{\raggedright\arraybackslash}p{(\linewidth - 12\tabcolsep) * \real{0.1429}}
  >{\raggedright\arraybackslash}p{(\linewidth - 12\tabcolsep) * \real{0.1429}}
  >{\raggedright\arraybackslash}p{(\linewidth - 12\tabcolsep) * \real{0.1429}}
  >{\raggedright\arraybackslash}p{(\linewidth - 12\tabcolsep) * \real{0.1429}}@{}}
\toprule\noalign{}
\begin{minipage}[b]{\linewidth}\raggedright
Target platform
\end{minipage} & \begin{minipage}[b]{\linewidth}\raggedright
Algorithms, Applications, Systems
\end{minipage} & \begin{minipage}[b]{\linewidth}\raggedright
Architecture, Networking
\end{minipage} & \begin{minipage}[b]{\linewidth}\raggedright
Scheduling
\end{minipage} & \begin{minipage}[b]{\linewidth}\raggedright
Green Computing
\end{minipage} & \begin{minipage}[b]{\linewidth}\raggedright
Modeling, Simulation
\end{minipage} & \begin{minipage}[b]{\linewidth}\raggedright
Other
\end{minipage} \\
\midrule\noalign{}
\endhead
\bottomrule\noalign{}
\endlastfoot
Cluster & 25 & 30 & 25 & 8 & 12 & 10 \\
Distributed & 10 & 2 & 10 & 5 & 10 & 5 \\
Cloud & 5 & 2 & 20 & 5 & 5 & 5 \\
IoT & 8 & 5 & 5 & 3 & 3 & 2 \\
Multi-Core/GPU & 0 & 0 & 3 & 0 & 2 & 1 \\
\end{longtable}
}

Figure 1: Publication counts by PDC domains and platforms for 176
research publications between 2016 and 2022 that include SimGrid-driven
simulation results. The sum of the counts is larger than 176 because
some publications target multiple domains and/or platform scenarios.

possible to design and implement a simulation framework that achieves
high accuracy, scalability, and versatility, as shown in the reference
SimGrid paper published in 2014 {[}2{]} and in many validation studies
{[}3, 4, 5, 6, 7, 8, 9, 10, 11, 12, 13, 14, 15{]}.

In {[}2{]} it is claimed that SimGrid could be used to develop
simulators for a broad range of PDC domains, thus achieving versatility.
We keep track of research publications that include simulation results
obtained with the SimGrid software {[}16{]}, which makes it possible to
verify the validity of that claim. Figure 1 shows the counts of the 176
research publications between 2016 and 2022 that use SimGrid (out of 646
such publications to date since 2001). We have categorized each
publication by target PDC domain and target platform scenario based on
our own inspection of each publication's content. The ``Other'' category
for the target PDC domains includes publications that focus on
fault-tolerance, model-checking, security, and education. Although our
categorization of some research publications may not be 100\% accurate,
the counts clearly show SimGrid's versatility.

Despite its popularity, SimGrid had severe shortcomings in terms of
extensibility and usability. First, it was difficult for users to
customize, extend, let alone replace, simulation models. While these
simulation models are sensible for the common case, the range of
research questions that users seek to answer using simulation is large
and often requires going beyond SimGrid's extant models. Second, SimGrid
only provided relatively low-level simulation abstractions and
corresponding APIs. These shortcomings were seen plainly in
user-provided feedback and through direct interactions with users during
SimGrid-focused events and hackathons, and also mentioned in the
literature {[}17{]}. Then, implementing a simulator of custom and/or
complex systems often required large software engineering and
development efforts.

In the last decade, since the publication of {[}2{]}, which used SimGrid
v3.10 (which we call v3 for simplicity), numerous efforts have focused
on improving SimGrid's extensibility and usability. The overall goal was
to lower or remove entry barriers for users to develop custom and/or
complex simulators, while retaining all of SimGrid's accuracy,
scalability, and versatility advantages. To achieve this goal, many
features have been released over a 10-year development cycle up to
v3.36, culminating in the SimGrid v4 release. While SimGrid has been the
subject of several previous publications, this work describes technical
developments and report on measures of success that have never been
previously published. Specifically:

\begin{enumerate}
\def\labelenumi{\arabic{enumi}.}
\tightlist
\item
  A description of the fundamental simulation abstractions in SimGrid v4
  (Section 3);\\
\item
  A description of mechanisms for customizing or replacing simulation
  models underlying these abstractions (Section 4);\\
\item
  A description of how these abstractions can be combined into
  higher-level abstractions (Section 5);\\
\item
  A demonstration of how a single simulator can employ several
  programming models (Section 5);\\
\item
  A survey of how the above has enabled a wide range of use cases
  (Section 6); and\\
\item
  A quantitative comparison of simulation scalability of SimGrid v3 and
  SimGrid v4 (Section 7).
\end{enumerate}

\section{2 Related work}\label{related-work}

There is a large literature devoted to the simulation of PDC platforms
and applications with many proposed PDC simulation frameworks. Some of
these frameworks have garnered sizable user communities and are still
actively maintained at the time of writing. They can be placed into two
broad categories based on the level of details of their underlying
simulation models, with different implications on accuracy, scalability,
and versatility.

Perhaps the most natural approach is to implement simulation models at a
high level of details, in an attempt to reproduce near-exact real-world
behaviors so as to achieve high accuracy. Many simulation frameworks
have followed this approach for simulating network resources
(packet-level simulators {[}18, 19{]}), compute resources
(cycle-accurate simulators {[}20, 21{]}), and I/O resources
(block-accurate simulators {[}22{]}). These frameworks are not versatile
from a PDC standpoint because they each focus on a particular class of
hardware components, and as such are typically used by researchers who
specialize in studying these components. While it is conceivable to
combine multiple such frameworks to create a versatile PDC simulation
framework, doing so is typically impractical due to scalability
limitations as simulation time is typically orders of magnitude longer
than simulated time. Some PDC simulation frameworks have been developed
that simulate net-

work communications at the packet level for high accuracy, with
typically a severe loss in scalability, but simulate other components at
lower levels of details {[}23, 24{]}.

One solution to address the scalability issue posed by simulation models
that implement a high level of detail is to employ Parallel Discrete
Event Simulation (PDES) {[}25{]}. This approach has been used
successfully by PDC simulation frameworks developed for simulating
high-performance computing (HPC) applications and platforms {[}26, 27,
28, 29, 30{]}. The use of PDES allows these frameworks to scale up the
platform on which the simulation is executed, possibly requiring a
platform of the same scale as that of the platform being simulated. By
contrast, many other PDC simulation frameworks, as discussed hereafter,
aim to achieve simulation scalability when executing the simulation on a
single compute node (or even a single core).

An alternate approach for achieving scalable simulations, without
scaling up compute resources, is to implement simulation models at a
relatively low level of detail. That is, rather than simulating
``microscopic'' behaviors of the target system, these models instead
rely on mathematical models that aim to capture ``macroscopic''
behaviors. For instance, rather than simulating the lifecycle of
individual network packets, they compute instantaneous data transfer
rates based on network path bandwidths and latencies and on current
network usage. The main challenge is to come up with such coarse-grain
simulation models that are reasonably accurate in spite of this low
level of detail {[}4{]}.

Several PDC frameworks have been developed using the above approach. The
most notable two such frameworks are SimGrid and GridSim {[}31{]}, both
of which have garnered large user communities, but using different
approaches for and achieving different levels of versatility, accuracy,
scalability, usability, and extensibility. SimGrid aims to be directly
versatile across a large range of PDC domains (see Figure 1). Instead,
versatility has been achieved in GridSim by implementing domain-specific
frameworks on top of it such as CloudSim {[}32{]} (which was later
re-implemented standalone reusing a subset of the GridSim simulation
abstractions) for cloud simulations, iFogSim {[}33{]} for IoT
simulations, DISSECT-CF {[}34{]} for IoT and cloud simulations, GroudSim
{[}35{]} for grid and cloud simulations, or OpenDC {[}36{]} and DCSim
{[}37{]} for data center simulations. The main focus of SimGrid up to v3
was the development of accurate and scalable simulation models via
extensive validation and invalidation studies {[}3, 4, 5, 6, 7, 8, 9,
10, 11, 12, 13, 14, 15{]}.

GridSim and the frameworks built on top of it have provided high
usability (due to providing higher-level simulation abstractions) and
high extensibility (as seen in the number of works that have
successfully extended CloudSim, e.g., {[}38, 39, 40, 41{]}). By
contrast, SimGrid v3 has had usability and extensibility shortcomings,
as already discussed in Section 1. This paper focuses on the technical
developments that have addressed these shortcomings, while preserving
simulation accuracy and scalability.

\section{3 Fundamental simulation
abstractions}\label{fundamental-simulation-abstractions}

SimGrid v4 provides convenient APIs, and efficient implementations
thereof, for three fundamental simulation abstractions that are
sufficient to describe and simulate virtually any PDC scenario:
resource, activity, and actor. SimGrid v3 provided abstractions similar
in intent, but suffered from a lack of separation of concern between the
models that determine the behaviors of resources, the activities that
consume these resources and the actors that launch these activities. The
refactoring and re-implementation of SimGrid in C++ in the last 10
years, was the occasion to enforce a strong separation between the user
and kernel spaces: each abstraction exposed to the user has its
counterpart in the simulation kernel. This design has been instrumental
for improving both extensibility and usability: the fundamental
abstractions can be customized and/or replaced in various ways (see
Section 4) and they can serve as the basis for higher-level abstractions
(see Section 5).

\section{3.1 Resources}\label{resources}

The first step in most SimGrid simulations is to describe a hardware
platform to simulate. This platform is composed of three types of basic
hardware resources based on the following three abstractions: (i) CPUs,
defined by a number of cores and a core compute speed; (ii) Network
links, defined by a latency and a bandwidth; and (iii) Disks, defined by
read and write bandwidths.

These basic resources are then aggregated and structured in higher-level
abstractions to describe a hardware platform either programmatically in
the simulator's code or in an external XML file. A physical host
comprises a CPU, any number of disks, and a network endpoint. Sim-Grid
provides also a virtual machine abstraction, which is a non-physical
host that can be instantiated on and moved between physical hosts. A
route is a vector of network links that defines a possibly multi-hop
network path between two network endpoints (e.g., between two hosts). A
route is either defined by the user or computed using some algorithm,
and the set of routes defines the physical network topology. A network
zone represents a set of hosts and links, has a network endpoint, and
knows how to determine the routes, if they exist, between two hosts in
the set and between a host and the zone's endpoint. For instance, a
network zone can be defined to represent a compute cluster in which all
nodes can communicate with each other and communicate with the outside
world via a shared endpoint. Finally, a platform is a hierarchy of
network zones. Determining the route between two hosts in differing
network zones is done recursively {[}42{]}.

\section{3.2 Activities}\label{activities}

\section{3.2.1 Activity abstraction}\label{activity-abstraction}

Given a hardware platform description, the goal of any simulator is to
simulate the execution of activities on the platform's hardware
resources: computations on CPUs, communications on network links, and
I/O operations on disks. SimGrid thus provides an activity abstraction,
and its API allows the user to write the code to create and manage the
execution of activities. All activities, regardless of on which
resources they are meant to run, follow the same lifecycle depicted in
Figure 2.

\textit{[Image omitted: images/b96022ca637a7a94c18d8c22786242e45f5c3db3aab3bbbb20392e7f6933cbab.jpg]}

flowchart

\begin{Shaded}
\begin{Highlighting}[]
\NormalTok{graph TD}
\NormalTok{    A["Init"] {-}{-}\textgreater{}|Start| B["Vetoed"]}
\NormalTok{    B {-}{-}\textgreater{}|Check| C["Running"]}
\NormalTok{    C {-}{-}\textgreater{} D["Cancel"]}
\NormalTok{    C {-}{-}\textgreater{} E["Fail"]}
\NormalTok{    C {-}{-}\textgreater{} F["Timeout"]}
\NormalTok{    C {-}{-}\textgreater{} G["Done"]}
\NormalTok{    C {-}{-}\textgreater{} H["Suspended"]}
\NormalTok{    H {-}{-}\textgreater{} I["Setter"]}
\end{Highlighting}
\end{Shaded}

Figure 2: SimGrid activity lifecycle.

Once created (Init state), configuration parameters or properties of an
activity can be set using setters, including setting dependencies on
other activities. The activity is then started, at which point it may
run immediately (Running state) or be held back (Vetoed state). The
activity may be in the Vetoed state because it depends on the completion
of other ongoing activities or because the resource on which the
activity must execute is not yet available or identified. While running,
an activity may be paused (Suspended state) and resumed any number of
times. When an activity terminates it can end up in four different
states depending on whether it terminates successfully (Done state), is
canceled by the user code (Cancel state), terminates unsuccessfully due
to a simulated resource failure or shutdown (Fail state), or reaches a
user-specified timeout (Timeout state).

\section{3.2.2 Activity simulation}\label{activity-simulation}

At the core of SimGrid are simulation models that determine the
completion date, in simulated time, of each activity. Each activity is
defined by a total amount of work to do (e.g., bytes to read, compute
operations to perform), an amount of work that remains to be done, and a
set of resources that are used to perform this work. The objective is to
compute the activity's completion date based on a notion of latency, in
seconds, and of a maximum rate of progress, or speed, in units of work
per time unit, both of which are computed based on the hardware
specifications of these resources. The latency is computed as function
(e.g., the sum) of the latencies of the resources. The speed is based on
a determination of the bottleneck resource and on the specifications of
that resource.

\textit{[Image omitted: images/7c724e0146d017ace5585414c0832c96a746a6798c75a7d0ec2130961e181410.jpg]}

flowchart

\begin{Shaded}
\begin{Highlighting}[]
\NormalTok{graph TD}
\NormalTok{    A["Actors"] {-}{-}\textgreater{} B["Activities"]}
\NormalTok{    B {-}{-}\textgreater{} C["Constraints"]}
\NormalTok{    C {-}{-}\textgreater{} D["Simulation kernel"]}
\NormalTok{    D {-}{-}\textgreater{} E["Linear MaxMin solver"]}
\NormalTok{    E {-}{-}\textgreater{} F["Simulation"]}
\NormalTok{    F {-}{-}\textgreater{} G["Host A"]}
\NormalTok{    F {-}{-}\textgreater{} H["Host B"]}
\NormalTok{    F {-}{-}\textgreater{} I["Host C"]}
\NormalTok{    F {-}{-}\textgreater{} J["Host D"]}
\NormalTok{    style A fill:\#f9f,stroke:\#333}
\NormalTok{    style B fill:\#ccf,stroke:\#333}
\NormalTok{    style C fill:\#cfc,stroke:\#333}
\NormalTok{    style D fill:\#fcc,stroke:\#333}
\NormalTok{    style E fill:\#cff,stroke:\#333}
\NormalTok{    style F fill:\#ffc,stroke:\#333}
\NormalTok{    style G fill:\#fcc,stroke:\#333}
\NormalTok{    style H fill:\#fcc,stroke:\#333}
\NormalTok{    style I fill:\#fcc,stroke:\#333}
\NormalTok{    style J fill:\#fcc,stroke:\#333}
\end{Highlighting}
\end{Shaded}

Figure 3: Overview of a SimGrid simulation.

At some point during the simulation multiple concurrent activities can
contend for some resources. In this case, resource shares must be
computed and allocated to each of these activities. SimGrid computes
these resource shares using a linear max-min (LMM) solver, which is
invoked at each simulation step. It solves a constrained optimization
problem to compute the speed of each activity. The default objective
function is to maximize the minimum speed across all activities. The
constraints impose that linear combinations of one or more activity
speeds are bounded by constants. These constants depend on the hardware
characteristics of the resources and on user-imposed constraints on the
rate of progress of activities. The coefficients in the linear
combinations are used to model various effects (see Section 4.3). Once
the speed of each activity has been computed, the time until the first
activity termination is computed, the simulation clock is advanced
accordingly, and the amounts of remaining work for all activities are
updated based on their speed and time elapsed. Completed Activities are
then removed from consideration for the next invocation of the solver.

Figure 3 depicts a 4-host platform with a 5-link ``dogbone'' network
topology (right-hand side). 4 concurrent activities are simulated, each
shown in a different color: two computation activities on Host A, a
communication activity that uses links B, C, and D, and another
communication that uses links C and E. The two computations contend for
Host A's compute capacity, and the two communications contend for Link
C's bandwidth. All these activities are defined in the simulation kernel
by total and remaining amounts of work. Each activity is also associated
to a variable \(\varrho\) , which is the activity's speed and whose
value needs to be determined by the LMM solver (left-hand side). The
constraints of the optimization problem are shown in the lower-left part
of the figure. There is one constraint for each resource used by ongoing
activities, where the right-hand side is the resource's capacity and the
left-hand side is a linear combination of the activity speeds, each
scaled by some factor (denoted as \(a_{x,y}\) where x is the resource
and y is the activity). For instance, the speed of the two communication
activities ( \(\varrho_{2}\) and \(\varrho_{3}\) ) both appear in the
constraint for Link C since both activities use some of that link's
bandwidth.

\textit{[Image omitted: images/12f066aa103ef56fd537837d4df233e42697a0244fe645d5a71c4d080a0dedd7.jpg]}

flowchart

\begin{Shaded}
\begin{Highlighting}[]
\NormalTok{graph TD}
\NormalTok{    A["Actor 1"] {-}{-}\textgreater{} B["Actor 2"]}
\NormalTok{    B {-}{-}\textgreater{} C["Actor 3"]}
\NormalTok{    C {-}{-}\textgreater{} D["Maestro"]}
\NormalTok{    D {-}{-}\textgreater{} E["Actor 2"]}
\NormalTok{    E {-}{-}\textgreater{} F["Actor 3"]}
\NormalTok{    F {-}{-}\textgreater{} G["Maestro"]}
\NormalTok{    style A fill:\#f9f,stroke:\#333}
\NormalTok{    style B fill:\#f9f,stroke:\#333}
\NormalTok{    style C fill:\#f9f,stroke:\#333}
\NormalTok{    style D fill:\#f9f,stroke:\#333}
\NormalTok{    style E fill:\#f9f,stroke:\#333}
\NormalTok{    style F fill:\#f9f,stroke:\#333}
\NormalTok{    style G fill:\#f9f,stroke:\#333}
\end{Highlighting}
\end{Shaded}

\begin{enumerate}
\def\labelenumi{(\alph{enumi})}
\tightlist
\item
  Simulation execution timeline in simulated time.
\end{enumerate}

\textit{[Image omitted: images/8827db941dd9289142d7ed263d3986e26efbee398e4cfd961810ebf639fc0495.jpg]}

flowchart

Simcall control flowchart showing interactions between actor, return,
handling, and time advances across three actors

\begin{enumerate}
\def\labelenumi{(\alph{enumi})}
\setcounter{enumi}{1}
\tightlist
\item
  Simulation execution timeline in wallclock time.\\
  Figure 4: Example simulation execution timeline for three actors
  between simulated times T1 and T2. (a) Simulated time -- Initially,
  the maestro is involved in running the LMM solver to update all
  activities' remaining amounts of work. These activities correspond to
  simcall invocations by the three actors, which, in this example, all
  complete at the same simulated time (T1). At this point, a scheduling
  round begins and each actor executes its code until the next simcall
  is placed (shown as blue line segments). Once all actors are blocks on
  a simcall, the maestro regains control and handles all these simcalls
  (B). In this example, the simcalls placed by Actor 2 and Actor 3 take
  zero simulated time. A new scheduling rounds begins and the code of
  these two actors is resumed (C) until they place a new simcall and
  become blocked again (D). Since all actors are blocked on simcalls,
  the maestro regains control, processes these two new simcalls, which,
  in this example, take non-zero simulated time. The LMM solver is
  invoked to update all activities' remaining amounts of work, thus
  advancing the simulated time until the first simcall completion (T2).
  This occurs for Actor 1, whose code is resumed (E). This process
  continues until all actors have completed, or until no actor can make
  further progress (which denotes a deadlock bug in the simulator). (b)
  Wallclock time -- The flow of control in the simulation's actual
  execution alternates between the maestro and the actors. Importantly,
  the code of two actors (in between the simcalls they place) never
  executes concurrently.
\end{enumerate}

See {[}2{]} for all details regarding SimGrid's simulation models, and
{[}3, 4, 5, 6, 7, 8, 9, 10, 11, 12, 13, 14, 15{]} for their experimental
validations.

\section{3.3 Actors}\label{actors}

\section{3.3.1 Actor abstraction}\label{actor-abstraction}

It is possible to write a SimGrid simulator that merely creates
activities (typically with dependencies), assigns them to resources, and
then launches the simulation, which ends when all activities have
completed. Most users, however, wish to simulate complex systems with
dynamic runtime behaviors using a more general Communicating Sequential
Processes (CSP) model. To do so, a SimGrid simulator can create one or
more actors. An actor is a (simulated) sequential process, defined by a
main procedure written by the user. This procedure can create activities
and manage their lifecycle. An actor can start an activity in three
modes: (i) Blocking -- the start call is blocking, and the corresponding
actor can only proceed when the activity terminates; (ii) Asynchronous
-- the start call returns immediately with a handle for the activity,
which provides an API to check for and to wait on the activity's
termination; (iii) Detached -- the start call returns immediately
without a handle in a ``fire and forget'' fashion. A simulation
typically comprises many actors, each of which is responsible for one or
more activities.

Conceptually actors execute on hosts independently, just like processes
would in a real-world platform. Actors can synchronize with each other
using classical synchronization abstractions (i.e., mutex, condition
variable, semaphore, barrier). These abstractions are implemented
internally as zero-time activities. There is another synchronization
abstraction called a mailbox. A mailbox is a logical rendez-vous point
through which actors can exchange messages, similar to a URL on which
one could post and retrieve data. Mailboxes are only used to match the
put and get requests of actors. Each such matching results in a
communication activity, either within a host when the communicating
actors run on the same host, or over an end-to-end network path when the
communicating actors run on different hosts.

\section{3.3.2 Actor simulation}\label{actor-simulation}

For the sake of deterministic execution, the actors cannot modify their
environment directly: each modification of the activities or other
interaction with the simulation kernel is serialized through a central
component that processes them in a deterministic order. For the sake of
modularity and correctness, SimGrid is designed as an operating system:
the actors issue simcalls (akin to system calls) to a simulation kernel
(akin to an OS kernel) called the maestro. The maestro decides which
actors can proceed and which ones must wait. Simcalls can be either
immediate if they take no time in the simulation (e.g., spawning another
actor), or blocking if their completion must wait for future events
(e.g., mutex locks require the mutex to be unlocked by its owner;
communications wait for the network to have provided enough
communication bandwidth for the data transfer to have completed).

A SimGrid simulation proceeds as a sequence of scheduling rounds. At
each round, code in all actors that are not currently blocked on a
simcall gets executed. This is done by the maestro, which passes the
control flow to the code of each actor in sequence. The control flow is
returned

to the maestro when the actor blocks on its next simcall. Once all
actors have executed up until the point where they have placed a
simcall, all simcalls are handled by the maestro. If some simcalls are
immediate, another scheduling round starts immediately for those actors
that are ready to execute. Once all actors have placed a blocking
simcall, the maestro invokes the LMM solver, which computes the time at
which the first pending simcall terminates. The simulation time is then
advanced to that point, one or more actors are unblocked, and a new
scheduling round begins. Figure 4a depicts scheduling rounds for an
example with three actors. From the simulation's perspective the code
written in between each simcall by the user for each actor takes zero
time; only simcalls can take time. If the user wishes to simulate the
time taken to execute this code, then a compute activity must be created
and started (via a simcall).

By default, the maestro and the actors execute sequentially on a single
core of the processor that runs the simulator (a multi-core execution
mode is also available {[}43{]}). Figure 4b shows the same example as
that shown in Figure 4a, but with the timeline drawn according to
wallclock time instead of simulated time. The execution flow alternates
between the maestro and the actors, and the code of two actors (in
between the simcalls they place) never executes concurrently. As a
result, although actors are implemented as threads of control within the
same address space, they can share memory without any risk of race
condition. This greatly simplifies the design and implementation of a
SimGrid simulator: although the simulator conceptually implements a
distributed system, this implementation can safely use globally visible
data structures for actors to share information in zero simulated time.
For instance, consider a simulator of a distributed storage system in
which there is a central directory of file locations, but the overhead
of accessing/updating this directory is deemed negligible by the user.
In this case, instead of implementing the in-simulation management of
this directory (i.e, implementing the directory as an actor with which
other actors must communicate via messages), the directory can simply be
implemented as, say, a dictionary data structure.

\section{3.3.3 Model checking ability}\label{model-checking-ability}

As stated in the previous section, the design by which all interactions
of the actors with their environment are properly mediated via simcalls
and the maestro ensures that the simulated execution is deterministic.
As a result, simulated execution events are observable and the causality
between these events is well-defined. This makes it possible to replace
the classical performance-oriented simulation engine with another engine
that only models the causal ordering of observed events for the purpose
of software model checking (SMC) {[}44{]}. The goal is to efficiently
explore all possible outcomes of the simulated application's execution
so as to assess the correctness of the application's implementation
(i.e., discover possible deadlocks, livelocks, message mismatches, and
other logical bugs). SimGrid v4 now comes with a model checker that
performs SMC and uses formal techniques to reduce the explored search
space by exploiting symmetries and prune causally equivalent histories.
This model checker can be used to verify MPI applications {[}45{]} and
pthread applications {[}46{]}. More generally, the correctness of any
application whose execution is simulated using SimGrid can be evaluated
in this manner (which, incidentally, also helped ironing out bugs in
SimGrid itself).

\section{4 Extension mechanisms}\label{extension-mechanisms}

A key concern for a simulation framework is its extensibility. This is
because users often want to model (part of their) target systems in
specific ways that do not correspond to the common case. In SimGrid v3
users had to understand and modify internals to extend simulation
abstractions and/or models. We have since added various mechanisms so
that the simulation framework can be extended conveniently.

\section{4.1 Rich and unified
resources}\label{rich-and-unified-resources}

In SimGrid v3, to answer the needs of many users, it was possible to
model Dynamic Voltage Frequency Scaling (DVFS) for CPU resources, by
which each CPU resource naturally implements a notion of pstate, where a
given state is defined by a numeric id and a compute speed, and can be
associated to a wattage. It was also possible to model background load
on CPU resources by attaching a load profile to a host that specifies
dynamic changes in the nominal speed according to a trace file or to a
random generator. Finally, it was also possible to attach to a host a
state profile to specify whether the host is up or down, which makes it
possible to simulate arbitrary host failures and churn patterns.

The above features were clearly useful but limited to CPUs and hosts.
Since SimGrid v3.11 they have been extended to the other resources
(network links and disks), thus unifying resource abstractions. As a
result, users can now modify the behavior of any resource by using the
pstate, load profile, and state profile features in creative ways to
simulate a range of complex phenomena, the simulation of which would
have been labor-intensive (i.e., requiring modifying internals) with
SimGrid v3. For instance, the startup period of any resource can now be
modeled easily with an extra pstate in which the resource performs no
work for activities on that resource. An actor can be in charge of
changing the resource's pstate after a given lapse of time that
corresponds to the duration of the startup period. Another example is
the simulation of Wi-Fi networks (introduced in SimGrid v3.26), where
the bandwidth experienced by a station is a function of the
signal-to-noise ratio between the station and the access point. It turns
out that this phenomenon can be simulated by assigning particular pstate
values to network links. Section 5.2 describes

how this technique, combined with other extension mechanisms, makes it
possible to model the behavior of Wi-Fi networks.

\section{4.2 Plugins}\label{plugins}

In SimGrid v4 we have developed a generic plugin feature, \(^{9}\) by
which arbitrary callbacks can be attached to signals that are fired by
specific simulation events (e.g., activity starts and completions, actor
creations and terminations). Plugins can be configured through arbitrary
properties attached to any resource when the simulated platform is
instantiated. They can also extend any simulation abstraction with any
arbitrary objects. These objects can be used to store persistent state
for the callbacks to use throughout the simulation.

Plugins have allowed us to not only simplify SimGrid's internal
implementation but also to make SimGrid more extensible. For instance,
SimGrid v3 provided a model of the energy consumption of CPU resources.
The implementation of this model was spread throughout SimGrid's core,
so as to pass relevant parameters to the energy model and invoke it to
update the energy consumption values at the correct simulated times.
This implementation was difficult to maintain. Furthermore, a user
wanting to modify (let alone extend or replace) the energy model had to
become essentially a full-fledged SimGrid developer. The energy model
was re-implemented as a self-contained plugin with its own configuration
(which specifies the consumption of each pstate in Watts) and callbacks
for the relevant events related to CPU resources and their use by
activities. The implementation of the plugin is confined to a single
source file which contains only 350 lines of C++ code, that is easy for
users to copy and modify.

The SimGrid distribution now comes with several built-in plugins (e.g.,
network link load and energy modeling (v3.18), WiFi network load and
energy modeling (v3.26), host energy modeling (v3.11), host load
modeling (v3.22), file system (v3.18), batteries and solar panels
(v3.34), computer room air handling (v3.35)). Users can also easily
develop their own plugins from scratch. For instance, often users wish
to collect and output information regarding specific events that happen
throughout the simulated execution, typically for post-mortem analysis.
Say that the user wishes to output to the terminal a time-stamped trace
of actor termination events. They can develop such a tracing plugin in
only a few lines of C++, as shown in Figure 5.

\section{4.3 Advanced modeling
mechanisms}\label{advanced-modeling-mechanisms}

As explained in Section 3.2, SimGrid's simulation core is implemented as
a constrained optimization problem solver, the LMM solver. Using the
SimGrid v4 API, there are three ways for users to modify the LMM so as
to extend SimGrid's simulation models, as described hereafter.

First, it is possible to simulate changes of the speed at which
resources perform activities, given the current set of simulated
activities, by defining ad hoc correction factors applied to the
resources' latencies and speeds. Such factors were already present in
SimGrid v3, for the specific purpose of simulating various network
protocol effects (e.g., for modeling TCP overhead {[}4{]} or adaptive
protocols used by MPI implementations {[}3{]}). But these factors were
global, static, only for network links and communication activities, and
were all hard-coded in the LMM. In SimGrid v4 we have extended these
correction factors to all resource kinds. A callback function can be
attached to each resource, at creation time or later. This function
takes the size of an activity as parameter and applies user-specified
correction factors. This allows users to extend SimGrid's simulation
core to, for instance, simulate CPU affinities to study classical
unrelated-machines scheduling problems, account for protocol change
depending on resources (e.g., Infiniband RDMA between distinct nodes,
memory copy for intra-node communications, cudaMemcpy when GPUs are
involved) or simulate the heteroscedastic performance of mechanical
disks.

\begin{Shaded}
\begin{Highlighting}[]
\NormalTok{SIMGRID\_REGISTER\_PLUGIN}\OperatorTok{(}\NormalTok{tracer}\OperatorTok{,} \StringTok{"Tracer"}\OperatorTok{,} \OperatorTok{\&}\NormalTok{tracer\_init}\OperatorTok{)}
\CommentTok{// Callback to place whenever an actor terminates}
\AttributeTok{static} \DataTypeTok{void}\NormalTok{ trace\_exit}\OperatorTok{(}\NormalTok{simgrid}\OperatorTok{::}\NormalTok{s4u}\OperatorTok{::}\NormalTok{Actor}\OperatorTok{\&}\NormalTok{ actor}\OperatorTok{)} \OperatorTok{\{}
    \BuiltInTok{std::}\NormalTok{cout }\OperatorTok{\textless{}\textless{}}\NormalTok{ simgrid}\OperatorTok{::}\NormalTok{s4u}\OperatorTok{::}\NormalTok{Engine}\OperatorTok{::}\NormalTok{get\_clock}\OperatorTok{()} \OperatorTok{\textless{}\textless{}} \StringTok{","}\OperatorTok{;}
    \BuiltInTok{std::}\NormalTok{cout }\OperatorTok{\textless{}\textless{}}\NormalTok{ actor}\OperatorTok{{-}\textgreater{}}\NormalTok{get\_name}\OperatorTok{()} \OperatorTok{\textless{}\textless{}} \BuiltInTok{std::}\NormalTok{endl}\OperatorTok{;}
\OperatorTok{\}}
\CommentTok{// Plugin initialization function}
\DataTypeTok{void}\NormalTok{ tracer\_init}\OperatorTok{()} \OperatorTok{\{}
    \CommentTok{// Attaching the callback to the relevant event}
\NormalTok{    simgrid}\OperatorTok{::}\NormalTok{s4u}\OperatorTok{::}\NormalTok{Actor}\OperatorTok{::}\NormalTok{on\_exit}\OperatorTok{(\&}\NormalTok{trace\_exit}\OperatorTok{);}
\OperatorTok{\}} 
\end{Highlighting}
\end{Shaded}

Figure 5: A simple ``tracing'' plugin.

Second, it is now possible to modify the LMM to simulate the fact that
the performance delivered by a resource can degrade with contention. In
other words, although by default the LMM simulates a linear sharing of a
resource among contending activities, the user can specify arbitrary,
including non-linear, sharing policies. Specifically, a user can attach
an arbitrary function to a resource, which computes the current
resource's speed as a function of the number of concurrent activities
using that resource. This function could implement any analytical model,
or could reproduce specific discrete behaviors observed on particular
real-world resources. This function is invoked by the LMM each time new
resource shares must be computed for that resource.

Combining these two methods, introduced in SimGrid v3.29, correction
factors and non-linear sharing behaviors, makes it possible to simulate
resource behaviors that go beyond what the default LMM can do. For
instance, it has allowed us to faithfully simulate the execution of I/O
benchmarks on several types of mechanical and solid-state disk drives
{[}7{]}.

Third, users can specify arbitrary concurrency limits for any resource.
When the number of activities that try to use this resource exceeds that
limit, the extraneous activities are delayed and have to wait for
currently running activities to

complete. Specifying different concurrency limits makes it possible to
extend the use of SimGrid to many scenarios, such as, the simulation of
network throttling on a multiprocessor system-on-chip, or to study
scheduling algorithms under classical theoretical constraints (e.g.,
1-port network model, strictly serial executions on a CPU).

The above mechanisms allow users to vastly modify and extend the default
behavior of the LMM so as to simulate a wide range of real-world
phenomena. But SimGrid also offers the capability to bypass the LMM
altogether. For instance, SimGrid can be compiled to use the well-known
ns-3 packet-level network simulator {[}18{]} as a network model.
Packet-level simulation accounts for the movement of every network
packet involved in every communication. This higher level of detail,
compared to the LMM-based models that only recompute the respective
instantaneous speeds of the currently ongoing communications when a
communication starts or stops, expectedly comes with a much higher
simulation time.

Finally, users can extend SimGrid with arbitrary models by overriding
two methods: next\_occurring\_event(), which returns the date of the
next event that will occur according to the model, and
update\_state(delta), which updates the model's state by shifting the
date forward by delta seconds. SimGrid v4 uses this plugin mechanism to
simulate the behavior of resources that are not handled by SimGrid's
extant models, e.g., batteries or solar panels that provide power to
hosts. The battery plugin can thus interrupt the simulation when a
battery becomes depleted, which requires to turn off the hosts that use
that battery as their primary power source. An FMI (Functional Mock-up
Interface) plugin for SimGrid has also been developed that makes it
possible to run a SimGrid simulation while co-simulating any FMI model
{[}47{]}, such as ones built with OpenModelica. This plugin leverages
the above extension mechanism to enable such co-simulation, interrupting
the SimGrid simulation when an event occurs in the FMI model, and
allowing SimGrid actors to modify the parameters of the FMI model during
the simulation.

\section{5 Better usability via composite abstractions and programming
models}\label{better-usability-via-composite-abstractions-and-programming-models}

Using the abstractions in Section 3 and the extension mechanisms in
Section 4, we have implemented several high-level, composite
abstractions in SimGrid v4. We have also made it possible to combine
multiple programming models within the same SimGrid simulator. The goal
is to increase usability by making simulator implementation easier for a
broad range of use cases.

\section{5.1 Composite platforms}\label{composite-platforms}

Although the abstractions described in Section 3.1 make it possible to
specify any conceivable hardware platform, doing so for large platforms
is time-consuming or labor-intensive, and almost always error-prone. One
of the most common parallel computing platforms is a homogeneous cluster
in which hosts are interconnected via some network topology. Describing
a cluster using the XML format would entail specifying: (1) each host
individually with its own name and speed; (2) a private link connecting
a host to a single backbone switch define by its latency and bandwidth;
(3) the backbone switch itself with its own latency and bandwidth
values; and (4) all the routes between each pair of hosts. Figure 6
shows such an XML description of a 64-node homogeneous cluster.

\begin{Shaded}
\begin{Highlighting}[]
\NormalTok{\textless{}}\KeywordTok{platform} \OtherTok{version=}\StringTok{"4.1"}\NormalTok{\textgreater{}}
\NormalTok{    \textless{}}\KeywordTok{zone} \OtherTok{id=}\StringTok{"cluster"} \OtherTok{routing=}\StringTok{"Full"}\NormalTok{\textgreater{}}
    \CommentTok{\textless{}!{-}{-} Declare the 64 hosts {-}{-}\textgreater{}}
\NormalTok{    \textless{}}\KeywordTok{host} \OtherTok{id=}\StringTok{"node{-}0.cluster"} \OtherTok{speed=}\StringTok{"1Gf"}\NormalTok{/\textgreater{}}
\NormalTok{    \textless{}}\KeywordTok{host} \OtherTok{id=}\StringTok{"node{-}1.cluster"} \OtherTok{speed=}\StringTok{"1Gf"}\NormalTok{/\textgreater{}}
\NormalTok{    [...]}
\NormalTok{    \textless{}}\KeywordTok{host} \OtherTok{id=}\StringTok{"node{-}63.cluster"} \OtherTok{speed=}\StringTok{"1Gf"}\NormalTok{/\textgreater{}}
    \CommentTok{\textless{}!{-}{-} Declare the 64 private links {-}{-}\textgreater{}}
\NormalTok{    \textless{}}\KeywordTok{link} \OtherTok{id=}\StringTok{"link\_0"} \OtherTok{bandwidth=}\StringTok{"10Gbps"} \OtherTok{latency=}\StringTok{"10us"}\NormalTok{/\textgreater{}}
\NormalTok{    \textless{}}\KeywordTok{link} \OtherTok{id=}\StringTok{"link\_1"} \OtherTok{bandwidth=}\StringTok{"10Gbps"} \OtherTok{latency=}\StringTok{"10us"}\NormalTok{/\textgreater{}}
\NormalTok{    [...]}
\NormalTok{    \textless{}}\KeywordTok{link} \OtherTok{id=}\StringTok{"link\_63"} \OtherTok{bandwidth=}\StringTok{"10Gbps"} \OtherTok{latency=}\StringTok{"10us"}\NormalTok{/\textgreater{}}
    \CommentTok{\textless{}!{-}{-} Declare the backbone switch {-}{-}\textgreater{}}
\NormalTok{    \textless{}}\KeywordTok{link} \OtherTok{id=}\StringTok{"backbone"} \OtherTok{bandwidth=}\StringTok{"100Gbps"} \OtherTok{latency=}\StringTok{"100us"}\NormalTok{/\textgreater{}}
    \CommentTok{\textless{}!{-}{-} Declare the (64 x 63) / 2 routes {-}{-}\textgreater{}}
\NormalTok{    \textless{}}\KeywordTok{route} \OtherTok{src=}\StringTok{"node{-}0.cluster"} \OtherTok{dst=}\StringTok{"node{-}1.cluster"}\NormalTok{\textgreater{}}
\NormalTok{    \textless{}}\KeywordTok{link\_ctn} \OtherTok{id=}\StringTok{"link\_0"}\NormalTok{/\textgreater{}}
\NormalTok{    \textless{}}\KeywordTok{link\_ctn} \OtherTok{id=}\StringTok{"backbone"}\NormalTok{/\textgreater{}}
\NormalTok{    \textless{}}\KeywordTok{link\_ctn} \OtherTok{id=}\StringTok{"link\_1"}\NormalTok{/\textgreater{}}
\NormalTok{    \textless{}/}\KeywordTok{route}\NormalTok{\textgreater{}}
\NormalTok{    [...]}
\NormalTok{    \textless{}}\KeywordTok{route} \OtherTok{src=}\StringTok{"node{-}62.cluster"} \OtherTok{dst=}\StringTok{"node{-}63.cluster"}\NormalTok{\textgreater{}}
\NormalTok{    \textless{}}\KeywordTok{link\_ctn} \OtherTok{id=}\StringTok{"link\_62"}\NormalTok{/\textgreater{}}
\NormalTok{    \textless{}}\KeywordTok{link\_ctn} \OtherTok{id=}\StringTok{"backbone"}\NormalTok{/\textgreater{}}
\NormalTok{    \textless{}}\KeywordTok{link\_ctn} \OtherTok{id=}\StringTok{"link\_63"}\NormalTok{/\textgreater{}}
\NormalTok{    \textless{}/}\KeywordTok{route}\NormalTok{\textgreater{}}
\NormalTok{    \textless{}/}\KeywordTok{zone}\NormalTok{\textgreater{}}
\NormalTok{\textless{}/}\KeywordTok{platform}\NormalTok{\textgreater{} }
\end{Highlighting}
\end{Shaded}

Figure 6: A partial verbose XML description of a homogeneous cluster.

Unsurprisingly, many SimGrid users target such cluster platforms (see
Figure 1). To increase usability for these users SimGrid's XML format
was extended (even before SimGrid v3) with built-in netzone definitions
that can be re-used, customized, and combined at will.

Figure 7 describes the same platform as in Figure 6 but using the
built-in tag that drastically simplifies the code users have to write.
The underlying network topology and routing of the clusters are computed
automatically {[}42{]}, in a way that is more efficient than when

\begin{Shaded}
\begin{Highlighting}[]
\NormalTok{\textless{}}\KeywordTok{platform} \OtherTok{version=}\StringTok{"4.1"}\NormalTok{\textgreater{}}
    \CommentTok{\textless{}!{-}{-} Declare a 64{-}node cluster {-}{-}\textgreater{}}
\NormalTok{    \textless{}}\KeywordTok{cluster} \OtherTok{id=}\StringTok{"cluster"} \OtherTok{prefix=}\StringTok{"node{-}"} \OtherTok{radical=}\StringTok{"0{-}63"} \OtherTok{suffix=}\StringTok{".cluster"} \OtherTok{speed=}\StringTok{"1Gf"} \OtherTok{bw=}\StringTok{"10Gbps"} \OtherTok{lat=}\StringTok{"10us"} \OtherTok{bb\_bw=}\StringTok{"100Gbps"} \OtherTok{bb\_lat=}\StringTok{"100us"}\NormalTok{/\textgreater{}}
\NormalTok{\textless{}/}\KeywordTok{platform}\NormalTok{\textgreater{} }
\end{Highlighting}
\end{Shaded}

Figure 7: A compact XML description of a homogeneous cluster.

done as in Figure 6. Additional built-in descriptions of complex network
topologies commonly found in HPC clusters, such as multi-dimensional
torus, Fat-Tree, or Dragonfly topologies are made available to the
users.

Describing platforms in XML is convenient as users can produce human-
and machine-readable descriptions relatively easily. Unfortunately, it
has proven difficult to evolve \(^{1}\) SimGrid's XML format to include
new capabilities while maintaining backward compatibility. To illustrate
these limitations let us consider the Summit leadership-class
supercomputer from Oak Ridge National Laboratory. SimGrid's XML format
makes it possible to declare a cluster with the desired number of nodes
(i.e., approximately 4,600 compute nodes) and the desired topology
(i.e., 3-level Fat-Tree network topology with 18 director switches)
using the appropriate properties in the tag. However, \(^{11}\) this tag
only supports compute nodes defined as simple multicore processors. As a
result, while the \(^{14}\) tag is convenient for some users, it is too
limiting for others who then have to revert to an overly complex XML
platform description. For this reason we have introduced, \(^{17}\) in
SimGrid v3.28, the capability to describe platforms programmatically
directly in the code of the simulator. This provides users with a much
greater flexibility than with XML, allowing them to combine the
fundamental abstractions (CPU, network link, disk) in arbitrary fashion.
For instance, in the Summit use case, it is straightforward to implement
a programmatic description of each compute node as comprising 2 CPUs, 6
GPUs, 2 NICs, and 1 NVMe, all interconnected via a custom internal
network topology. Note that although SimGrid does not provide a built-in
GPU model, custom GPU models can be easily built by combining SimGrid's
fundamental abstractions {[}48{]}.

Figure 8 shows a programmatic C++ description of a more complex platform
with two independent hosts, representing a coordinator and a database
service, each connected to three homogeneous clusters of different
sizes. \(^{38}\) While this way of describing platforms may seem more
complicated at first, it is much easier to evolve than its XML
counterpart. For instance, adding three more clusters to that platform
to simulate a larger scale scenario would simply amount to adding three
more sizes in the vector (line 42), but would require adding three more
\(^{44}\) tags and declare all the additional routes in the XML
description, which is error-prone. Moreover, increasing the complexity
of the internal structure of the cluster can be done programmatically,
while it is highly constrained by the XML format, as explained earlier.
\(^{46}\)

\section{5.2 Composite network routes}\label{composite-network-routes}

By default, network topologies specified in simulated platforms assume
classical TCP-based end-to-end network routes. But expressing composite
routes where different behaviors are simulated for different portions of
the routes can be compelling. This is the case of the simulation of Wi-

\begin{Shaded}
\begin{Highlighting}[]
\NormalTok{NetZone}\OperatorTok{*}\NormalTok{ create\_cluster}\OperatorTok{(}\NormalTok{NetZone}\OperatorTok{*}\NormalTok{ root}\OperatorTok{,} \BuiltInTok{std::}\NormalTok{string suffix}\OperatorTok{,} \DataTypeTok{int}\NormalTok{ num\_hosts}\OperatorTok{)} \OperatorTok{\{}
    \KeywordTok{auto}\OperatorTok{*}\NormalTok{ cluster }\OperatorTok{=}\NormalTok{ create\_star\_zone}\OperatorTok{(}\StringTok{"cluster"}\OperatorTok{+}\NormalTok{suffix}\OperatorTok{){-}\textgreater{}}\NormalTok{set\_parent}\OperatorTok{(}\NormalTok{root}\OperatorTok{);}

    \CommentTok{// create gateway}
\NormalTok{    cluster}\OperatorTok{{-}\textgreater{}}\NormalTok{set\_gateway}\OperatorTok{(}\NormalTok{cluster}\OperatorTok{{-}\textgreater{}}\NormalTok{create\_router}\OperatorTok{(}\StringTok{"cluster"} \OperatorTok{+}\NormalTok{ suffix }\OperatorTok{+} \StringTok{"{-}router"}\OperatorTok{));}

    \CommentTok{// create the backbone link}
    \KeywordTok{auto}\OperatorTok{*}\NormalTok{ backbone }\OperatorTok{=}\NormalTok{ cluster}\OperatorTok{{-}\textgreater{}}\NormalTok{create\_link}\OperatorTok{(}\StringTok{"backbone"} \OperatorTok{+}\NormalTok{ suffix}\OperatorTok{,} \StringTok{"100Gbps"}\OperatorTok{){-}\textgreater{}}\NormalTok{set\_latency}\OperatorTok{(}\StringTok{"100us"}\OperatorTok{);}

    \CommentTok{// create all hosts and connect them to outside world}
    \ControlFlowTok{for} \OperatorTok{(}\DataTypeTok{int}\NormalTok{ i }\OperatorTok{=} \DecValTok{0}\OperatorTok{;}\NormalTok{ i }\OperatorTok{\textless{}}\NormalTok{ num\_hosts}\OperatorTok{;}\NormalTok{ i}\OperatorTok{++)} \OperatorTok{\{}
    \BuiltInTok{std::}\NormalTok{string name }\OperatorTok{=} \StringTok{"node{-}"} \OperatorTok{+} \BuiltInTok{std::}\NormalTok{to\_string}\OperatorTok{(}\NormalTok{i}\OperatorTok{)+}\NormalTok{suffix}\OperatorTok{;}
    \CommentTok{// create host}
    \KeywordTok{auto}\OperatorTok{*}\NormalTok{ host }\OperatorTok{=}\NormalTok{ cluster}\OperatorTok{{-}\textgreater{}}\NormalTok{create\_host}\OperatorTok{(}\NormalTok{name}\OperatorTok{,} \StringTok{"1Gf"}\OperatorTok{);}
    \CommentTok{// create link}
    \KeywordTok{auto}\OperatorTok{*}\NormalTok{ link }\OperatorTok{=}\NormalTok{ cluster}\OperatorTok{{-}\textgreater{}}\NormalTok{create\_link}\OperatorTok{(}\NormalTok{name}\OperatorTok{+}\StringTok{"\_link"}\OperatorTok{,} \StringTok{"10Gbps"}\OperatorTok{){-}\textgreater{}}\NormalTok{set\_latency}\OperatorTok{(}\StringTok{"10us"}\OperatorTok{);}
    \CommentTok{// add route between host and any other host}
\NormalTok{    cluster}\OperatorTok{{-}\textgreater{}}\NormalTok{add\_route}\OperatorTok{(}\NormalTok{host}\OperatorTok{,} \KeywordTok{nullptr}\OperatorTok{,} \OperatorTok{\{}\NormalTok{link}\OperatorTok{,}\NormalTok{ backbone}\OperatorTok{\});}
    \OperatorTok{\}}
\NormalTok{    cluster}\OperatorTok{{-}\textgreater{}}\NormalTok{seal}\OperatorTok{();}
    \ControlFlowTok{return}\NormalTok{ cluster}\OperatorTok{;}
\OperatorTok{\}}

\DataTypeTok{int}\NormalTok{ main}\OperatorTok{(}\DataTypeTok{int}\NormalTok{ argc}\OperatorTok{,} \DataTypeTok{char}\OperatorTok{**}\NormalTok{ argv}\OperatorTok{)} \OperatorTok{\{}
    \CommentTok{// Create the platform}
    \KeywordTok{auto}\OperatorTok{*}\NormalTok{ root }\OperatorTok{=}\NormalTok{ create\_full\_zone}\OperatorTok{(}\StringTok{"world"}\OperatorTok{);}

    \CommentTok{// Create a coordinator zone/host}
    \KeywordTok{auto}\NormalTok{ coordinator\_zone }\OperatorTok{=}\NormalTok{ create\_full\_zone}\OperatorTok{(}\StringTok{"Coordinator"}\OperatorTok{){-}\textgreater{}}\NormalTok{set\_parent}\OperatorTok{(}\NormalTok{root}\OperatorTok{);}
\NormalTok{    coordinator\_zone}\OperatorTok{{-}\textgreater{}}\NormalTok{create\_host}\OperatorTok{(}\StringTok{"coordinator.org"}\OperatorTok{,} \StringTok{"1Gf"}\OperatorTok{);}
\NormalTok{    coordinator\_zone}\OperatorTok{{-}\textgreater{}}\NormalTok{seal}\OperatorTok{();}

    \CommentTok{// Create a database zone/host}
    \KeywordTok{auto}\NormalTok{ database\_zone }\OperatorTok{=}\NormalTok{ create\_full\_zone}\OperatorTok{(}\StringTok{"Database"}\OperatorTok{){-}\textgreater{}}\NormalTok{set\_parent}\OperatorTok{(}\NormalTok{root}\OperatorTok{);}
\NormalTok{    database\_zone}\OperatorTok{{-}\textgreater{}}\NormalTok{create\_host}\OperatorTok{(}\StringTok{"database.org"}\OperatorTok{,} \StringTok{"1Gf"}\OperatorTok{){-}\textgreater{}}\NormalTok{create\_disk}\OperatorTok{(}\StringTok{"db"}\OperatorTok{,} \StringTok{"100MBps"}\OperatorTok{,} \StringTok{"50MBps"}\OperatorTok{);}
\NormalTok{    database\_zone}\OperatorTok{{-}\textgreater{}}\NormalTok{seal}\OperatorTok{();}

    \CommentTok{// Create a single link as a simple abstraction of the whole wide{-}area network}
    \KeywordTok{auto}\OperatorTok{*}\NormalTok{ internet }\OperatorTok{=}\NormalTok{ root}\OperatorTok{{-}\textgreater{}}\NormalTok{create\_link}\OperatorTok{(}\StringTok{"internet"}\OperatorTok{,} \StringTok{"200MBps"}\OperatorTok{){-}\textgreater{}}\NormalTok{set\_latency}\OperatorTok{(}\StringTok{"1ms"}\OperatorTok{);}

    \CommentTok{// Create three clusters and connect them to the coordinator and database}
    \BuiltInTok{std::}\NormalTok{vector}\OperatorTok{\textless{}}\DataTypeTok{int}\OperatorTok{\textgreater{}}\NormalTok{ cluster\_sizes }\OperatorTok{=} \OperatorTok{\{}\DecValTok{16}\OperatorTok{,} \DecValTok{32}\OperatorTok{,} \DecValTok{40}\OperatorTok{\};}
    \DataTypeTok{int}\NormalTok{ i }\OperatorTok{=} \DecValTok{0}\OperatorTok{;}
    \ControlFlowTok{for} \OperatorTok{(}\KeywordTok{auto}\NormalTok{ size }\OperatorTok{:}\NormalTok{ cluster\_sizes}\OperatorTok{)} \OperatorTok{\{}
    \KeywordTok{auto}\OperatorTok{*}\NormalTok{ cluster }\OperatorTok{=}\NormalTok{ create\_cluster}\OperatorTok{(}\NormalTok{root}\OperatorTok{,} \StringTok{".cluster"}\OperatorTok{+}\BuiltInTok{std::}\NormalTok{to\_string}\OperatorTok{(}\NormalTok{i}\OperatorTok{++)} \OperatorTok{+} \StringTok{".org"}\OperatorTok{,}\NormalTok{ size}\OperatorTok{);}
\NormalTok{    root}\OperatorTok{{-}\textgreater{}}\NormalTok{add\_route}\OperatorTok{(}\NormalTok{coordinator\_zone}\OperatorTok{,}\NormalTok{ cluster}\OperatorTok{,} \OperatorTok{\{}\NormalTok{internet}\OperatorTok{\});}
\NormalTok{    root}\OperatorTok{{-}\textgreater{}}\NormalTok{add\_route}\OperatorTok{(}\NormalTok{database\_zone}\OperatorTok{,}\NormalTok{ cluster}\OperatorTok{,} \OperatorTok{\{}\NormalTok{internet}\OperatorTok{\});}
    \OperatorTok{\}}
\NormalTok{    root}\OperatorTok{{-}\textgreater{}}\NormalTok{seal}\OperatorTok{();}
\OperatorTok{\}} 
\end{Highlighting}
\end{Shaded}

Figure 8: Programmatic C++ description a platform composed of two
individual hosts and three homogeneous clusters of different sizes.

Fi networks, for instance, needed to study IoT platforms. This requires
a new type of network zone that comprises the access points of the Wi-Fi
network, the hosts (stations in the Wi-Fi terminology) connected to it,
and a single network link declared with a specific bandwidth sharing
policy adapted to Wi-Fi networks. This pseudo-link is then included into
the network routes between any two stations within the Wi-Fi zone.

One fundamental difference between Wi-Fi and wired networks is that the
performance of the former is not determined by the bandwidths and
latencies of network links but by two characteristics of the access
point, defined as properties of the zone. First, the Modulation and
Coding Scheme (MCS) defines the speed at which the access point
exchanges data with all the stations. This speed directly depends on the
access point's model and configuration, and also on the distance between
the involved station and access point. More precisely, the data rate
depends on the signal-noise ratio (SNR) between the communicating
entities. Second, the number of antennas defines the amount of Spatial
Streams that the access point can simultaneously serve. In practice, a
given access point will provide several levels of performance (called
Data Rates) depending on its hardware characteristics. In SimGrid, the
Wi-Fi pseudo-link is given a set of pstates representing each of these
data rates (see Section 4.1). Each host connected to this access point
may have a different pstate on that link, representing the MCS that
would be used between the station and the access point corresponding to
the SNR between these elements. Additionally, the LMM is extended to
model non-linear resource sharing (see Section 4.3) so that the maximum
throughput of a Wi-Fi network is not a constant but rather a function of
the number of concurrent flows {[}49{]}. Implementing a fast but
accurate model of Wi-Fi performance would have been much more difficult
without these extension mechanisms provided since v3.26 (pstates for any
resource, user-provided resource sharing function).

\section{5.3 Composite activities}\label{composite-activities}

When writing a simulator for complex application execution scenarios,
actors often need to start multiple activities of different types
asynchronously and manage their concurrent execution. For instance, when
writing a simulator of an iterative numerical simulation based on domain
decomposition, each actor may have to exchange data with its neighbors,
write intermediate results to the file system, and perform computations,
all concurrently. To ease the management of concurrent asynchronous
activities, SimGrid v3.35 introduced an ActivitySet data container, in
which asynchronous activities of any types can be stored. This container
allows users to simply test or wait for the completion of any or all the
activities.

As seen in Figure 2, activities have a well-defined lifecycle that
corresponds to the unique execution of a certain amount of work.
However, in some execution models, such as Synchronous Data Flow (SDF)
or Bulk Synchronous Parallel (BSP), the same activity, or activity
sequence, must be repeated several times.

To support such execution models, SimGrid v3.34 introduced the concept
of tasks. A task has an underlying activity that can be repeated
multiple times. Tasks can be organized in graphs to simulate complex
iterative execution patterns with inter-task control dependencies.
SimGrid v4 also introduced the notion of tokens that circulate through a
task graph and can carry any user-defined data to implement inter-task
data dependencies. When all dependencies of a task become satisfied it
fires a new instance of its underlying communication, computation, or
I/O activity. It then blocks until all its dependencies become satisfied
again. The parameters and successors of a task can be redefined at
runtime using the same callback/signal mechanism as that used for
plugins. This allows users to dynamically change the topology of the
task graph, implement conditional branches, or introduce variability in
the duration of the simulated activities.

The task and token abstractions can be used to simulate relatively
simple, yet commonplace, designs which involve pipelined activities. As
an example, Figure 9 illustrates the usage of tasks and tokens to
implement an iterative ``computation and data exchange'' pattern between
two hosts, which is automatically repeated N times. That is, at each
iteration, each host performs a computation task and then a
communication task to simulate the exchange of data over the network.
This example also includes the simulation of a periodic checkpoint that
must occur every \(n \ll N\) iterations. This is easily simulated via a
token sent every n iterations, which triggers the execution of an I/O
task.

\textit{[Image omitted: images/d6300f145d11247ffeaf2bb6cdc5b871a94ad1e5cc634c84486ca8989d20ef90.jpg]}

flowchart

\begin{Shaded}
\begin{Highlighting}[]
\NormalTok{graph LR}
\NormalTok{    A["Host{-}1"] {-}{-}\textgreater{} B["Exec"]}
\NormalTok{    B {-}{-}\textgreater{} C["Comm"]}
\NormalTok{    D["Host{-}2"] {-}{-}\textgreater{} E["Exec"]}
\NormalTok{    E {-}{-}\textgreater{} F["Comm"]}
\NormalTok{    F {-}{-}\textgreater{} G["I/O"]}
\NormalTok{    C {-}{-}\textgreater{} G}
\NormalTok{    F {-}{-}\textgreater{} G}
\end{Highlighting}
\end{Shaded}

Figure 9: A simple BSP execution implemented with tasks and tokens. Two
hosts perform a computation and then exchange their data. This pattern
is automatically repeated N times. A conditional checkpoint is triggered
by the emission of a token sent every \(n \ll N\) iterations.

Although the above task and token abstractions are powerful, they can
become relatively costly in terms of simulation time due to each
operation being simulated as an individual task. This would be the case,
for instance, for simulating the buffered transfer of data stored on
disk to a remote host, which involves fine-grain pipelining of I/O
operations and communications. For this reason, SimGrid v4 provides
simulation abstractions that implement these types of concurrent
operations using a continuous fluid approximation. This approximation
does not simulate the discrete execution of these operations, but
instead computes the bottleneck operation and makes all operations
proceed at the speed of that bottleneck. The simula-

tion is thus implemented at a coarser grain, trading off simulation
details for faster simulation execution. Specifically, two such
abstractions are provided: parallel task, or ptask, and I/O stream
(introduced in v3.31). A ptask encapsulates computations and
communications to be executed over a set of interconnected hosts (e.g.,
an iterative matrix multiply parallel computation). It is defined by
volumes of computation to be executed at each host and volumes of data
to transfer between each host pair. An I/O stream encapsulates I/O
operations and communications for sending data between hosts and/or
disks (e.g., a video stream application with repeating I/O reads on a
host, each followed by network transfer to a remote host).

\section{5.4 Composite programming
models}\label{composite-programming-models}

The simulation abstractions presented in Section 3 allow users to
implement a simulator using the CSP programming model, which is general
and can be used to implement a simulator of virtually any distributed
system. However, SimGrid also supports other, less general, programming
models that are specialized for common use cases, so that simulators for
these use cases can be implemented with minimal effort.

As mentioned at the beginning of Section 3.3, it is possible to
implement a SimGrid simulator without creating any actor: the simulator
simply creates activities and sets up dependencies between these
activities. Since SimGrid v3.30, it is possible to describe a static
Directed Acyclic Graph (DAG) of activities. Each activity in this DAG
can be a computation, a network communication, an I/O operation, or any
of the composite activities described in Section 5.3. As a result, it is
possible to describe a static data-flow application, and the simulator
merely specifies on which resource(s) each activity is to be executed.
It can then launch the simulation until completion of all activities or
until a particular lapse of (simulated) time has elapsed. This
programming model is more restricted than the CSP model, but proves
useful and convenient for many users. Figure 1 shows that many SimGrid
users develop simulators for evaluating scheduling algorithms. Many of
these simulators are used to study applications structured as static
DAGs, for which different schedules are computed and must be evaluated
in simulation. The above programming model makes it straightforward to
implement such simulators.

A commonplace programming model for cluster platforms (the most
frequently targeted class of platform according to Figure 1) is
distributed-memory programming using message-passing. In practice, this
programming model is implemented using the Message Passing Interface
(MPI) standard {[}50{]}, which is not a fully general CSP programming
model but provides many convenient higher-level abstractions such as
collective communications. For this reason, SimGrid provides an
implementation of MPI, called SMPI (Simulated MPI). It makes it possible
for users to write standard MPI programs, or use unmodified existing MPI
programs, and execute them seamlessly in simulation. SMPI uses several
techniques to ensure that the simulated executions can be executed on a
single machine scalably {[}15{]}. It also implements accurate simulation
models of both point-to-point and MPI collective communication
operations, accounting for the specific schemes and algorithms
implemented in particular implementations of the MPI standard (which the
user can select at will). These models have been thoroughly validated,
and we refer the reader to {[}3, 9, 51{]} for experimental validation
results and all technical details.

A key usability enhancement in SimGrid v4 is that all the above
programming models can now be combined at will within a single
simulator: different components of a simulator can use different
programming models.

To illustrate this capability, we present a full-fledged example.
Consider a coordinator-worker application where each worker executes on
a node of a compute cluster. When idle, each worker requests work from a
coordinator that is running on some remote host. This coordinator-worker
scheme can be implemented naturally using SimGrid's CSP programming
model. Say that the work that each worker must perform consists in
invoking an MPI program that executes on all compute nodes of its
cluster. At each iteration of this MPI program the MPI ranks perform a
computation and then synchronize via an all-to-all communication. This
can be implemented easily, since v3.34, using SimGrid's MPI programming
model. Say now that, as part of this computation, at each iteration the
MPI process with rank 0 must upload 1MB of data to a remote database.
This consists in sending 1MB to a remote host and then writing 1MB to a
disk at that host. This can be easily implemented as a 2-activity
sequence using SimGrid's data-flow programming model. While the general
CSP programming model would be conceptually sufficient to implement the
entire simulator, the ability to combine multiple programming models
vastly reduces the overall development effort.

We have produced a working implementation of the above example using
SimGrid's C++ API in less than 150 lines of code (code available on
GitHub {[}52{]}). We show here relevant code fragments, all redacted for
brevity. Most of the code of the main() function (Figure 10) defines the
simulated hardware (as already shown in Figure 8). A few lines of code
are used to state that one Coordinator actor is to be started on some
host in the platform, and one Worker actor is to be started on a compute
node of each of 3 clusters (lines 9-13).

Figure 11 shows the Coordinator code which first creates its mailbox
(line 5), and then constructs a queue of workunits for workers to
perform (lines 7-9). Each workunit is described by a number of
iterations, an amount of data to communicate, and an amount of work to
compute. The workunit queue ends with 3 ``poison pills'' through which
workers will discover that there is no more work to be done. The
coordinator then goes through a loop

\begin{Shaded}
\begin{Highlighting}[]
\DataTypeTok{int}\NormalTok{ main}\OperatorTok{(}\DataTypeTok{int}\NormalTok{ argc}\OperatorTok{,} \DataTypeTok{char} \OperatorTok{**}\NormalTok{argv}\OperatorTok{)} \OperatorTok{\{}
    \CommentTok{// Create a simulation engine}
    \KeywordTok{auto}\NormalTok{ engine }\OperatorTok{=} \KeywordTok{new}\NormalTok{ Engine}\OperatorTok{(\&}\NormalTok{argc}\OperatorTok{,}\NormalTok{ argv}\OperatorTok{);}
    \CommentTok{// Create a simulated platform}
    \BuiltInTok{std::}\NormalTok{vector}\OperatorTok{\textless{}}\DataTypeTok{int}\OperatorTok{\textgreater{}}\NormalTok{ cluster\_sizes }\OperatorTok{=} \OperatorTok{\{}\DecValTok{16}\OperatorTok{,} \DecValTok{32}\OperatorTok{,} \DecValTok{40}\OperatorTok{\};}
    \OperatorTok{[...]}
    \CommentTok{// Create 1 coordinator}
\NormalTok{    Actor}\OperatorTok{::}\NormalTok{create}\OperatorTok{(}\StringTok{"Coordinator"}\OperatorTok{,}
\NormalTok{    Host}\OperatorTok{::}\NormalTok{by\_name}\OperatorTok{(}\StringTok{"coordinator.org"}\OperatorTok{),}\NormalTok{ Coordinator}\OperatorTok{());}
    \CommentTok{// Create 3 worker actors}
\NormalTok{    Actor}\OperatorTok{::}\NormalTok{create}\OperatorTok{(}\StringTok{"Worker1"}\OperatorTok{,}
\NormalTok{    Host}\OperatorTok{::}\NormalTok{by\_name}\OperatorTok{(}\StringTok{"host{-}0.cluster1.org"}\OperatorTok{),}\NormalTok{ Worker}\OperatorTok{());}
\NormalTok{    Actor}\OperatorTok{::}\NormalTok{create}\OperatorTok{(}\StringTok{"Worker2"}\OperatorTok{,}
\NormalTok{    Host}\OperatorTok{::}\NormalTok{by\_name}\OperatorTok{(}\StringTok{"host{-}0.cluster2.org"}\OperatorTok{),}\NormalTok{ Worker}\OperatorTok{());}
\NormalTok{    Actor}\OperatorTok{::}\NormalTok{create}\OperatorTok{(}\StringTok{"Worker3"}\OperatorTok{,}
\NormalTok{    Host}\OperatorTok{::}\NormalTok{by\_name}\OperatorTok{(}\StringTok{"host{-}0.cluster3.org"}\OperatorTok{),}\NormalTok{ Worker}\OperatorTok{());}
    \CommentTok{// Launch the simulation}
\NormalTok{    engine}\OperatorTok{{-}\textgreater{}}\NormalTok{run}\OperatorTok{();}
\OperatorTok{\}} 
\end{Highlighting}
\end{Shaded}

Figure 10: main() code for the example in Section 5.4.

(lines 15-22) in which it waits for workers to send it a mailbox to
which it replies with the next workunit to be performed (via a simulated
128-byte message), until all workunits have been performed.

\begin{Shaded}
\begin{Highlighting}[]
\KeywordTok{class}\NormalTok{ Coordinator }\OperatorTok{\{}
    \KeywordTok{public}\OperatorTok{:}
    \DataTypeTok{void} \KeywordTok{operator}\OperatorTok{()} \OperatorTok{()} \OperatorTok{\{}
    \CommentTok{// Create my mailbox}
    \KeywordTok{auto}\NormalTok{ my\_mailbox }\OperatorTok{=}\NormalTok{ Mailbox}\OperatorTok{::}\NormalTok{by\_name}\OperatorTok{(}\StringTok{"coordinator\_mb"}\OperatorTok{);}
    \CommentTok{// Create 10 workunits: iter=100, size=10MB, work=2Gf}
    \BuiltInTok{std::}\NormalTok{deque}\OperatorTok{\textless{}}\NormalTok{WorkUnit}\OperatorTok{*\textgreater{}}\NormalTok{ todo}\OperatorTok{;}
    \ControlFlowTok{for} \OperatorTok{(}\DataTypeTok{int}\NormalTok{ i }\OperatorTok{=} \DecValTok{0}\OperatorTok{;}\NormalTok{ i }\OperatorTok{\textless{}} \DecValTok{10}\OperatorTok{;}\NormalTok{ i}\OperatorTok{++)}
\NormalTok{    todo}\OperatorTok{.}\NormalTok{push\_front}\OperatorTok{(}\KeywordTok{new}\NormalTok{ WorkUnit}\OperatorTok{(}\DecValTok{100}\OperatorTok{,} \DecValTok{10}\OperatorTok{*}\NormalTok{MB}\OperatorTok{,} \DecValTok{2}\OperatorTok{*}\NormalTok{Gf}\OperatorTok{));}
    \CommentTok{// Add "poison pills" for worker terminations}
    \ControlFlowTok{for} \OperatorTok{(}\DataTypeTok{int}\NormalTok{ i }\OperatorTok{=} \DecValTok{0}\OperatorTok{;}\NormalTok{ i }\OperatorTok{\textless{}}\NormalTok{ num\_workers}\OperatorTok{;}\NormalTok{ i}\OperatorTok{++)}
\NormalTok{    todo}\OperatorTok{.}\NormalTok{push\_front}\OperatorTok{(}\KeywordTok{new}\NormalTok{ WorkUnit}\OperatorTok{(}\DecValTok{0}\OperatorTok{,} \DecValTok{0}\OperatorTok{,} \DecValTok{0}\OperatorTok{));}

    \CommentTok{// Main loop}
    \ControlFlowTok{while} \OperatorTok{(}\KeywordTok{not}\NormalTok{ todo}\OperatorTok{.}\NormalTok{empty}\OperatorTok{())} \OperatorTok{\{}
    \CommentTok{// Wait for a worker to send me their mailbox}
    \KeywordTok{auto}\NormalTok{ worker\_mailbox }\OperatorTok{=}\NormalTok{ my\_mailbox}\OperatorTok{{-}\textgreater{}}\NormalTok{get}\OperatorTok{\textless{}}\NormalTok{Mailbox}\OperatorTok{\textgreater{}();}
    \CommentTok{// Reply with the next workunit (128{-}byte message)}
\NormalTok{    worker\_mailbox}\OperatorTok{{-}\textgreater{}}\NormalTok{put }\OperatorTok{(}\NormalTok{todo}\OperatorTok{.}\NormalTok{back}\OperatorTok{(),} \DecValTok{128}\OperatorTok{);}
    \CommentTok{// Remove workunit from queue}
\NormalTok{    todo}\OperatorTok{.}\NormalTok{pop\_back}\OperatorTok{();}
    \OperatorTok{\}}
\OperatorTok{\}} 
\end{Highlighting}
\end{Shaded}

Figure 11: Coordinator C++ functor for the example in Section 5.4.

Figure 12 shows the Worker code. First, the worker identifies the
database's host and disk (lines 4-6) and the compute nodes in its
cluster (line 8), after which it creates its own mailbox (line 10) and
loops (line 13). At each iteration of the loop it requests work from the
coordinator via a 32-byte message (line 15), receives a workunit (line
17), and aborts if the workunit is a poison pill (line 19). If the
workunit is not a poison pill, then it asynchronously starts a simulated
MPI program at line 22. The first argument to the
SMPI\_app\_instance\_start() function is a name that will be used to
later wait for the completion of the MPI program. The second argument is
a lambda expression whose code is standard MPI code (lines 23-53). The
third argument is the list of compute nodes on which to execute the MPI
program (line 54).

\begin{Shaded}
\begin{Highlighting}[]
\KeywordTok{class}\NormalTok{ Worker }\OperatorTok{\{}
    \KeywordTok{public}\OperatorTok{:}
    \DataTypeTok{void} \KeywordTok{operator}\OperatorTok{()} \OperatorTok{()} \OperatorTok{\{}
    \CommentTok{// Get database host and disk}
    \KeywordTok{auto}\NormalTok{ dbhost }\OperatorTok{=}\NormalTok{ Host}\OperatorTok{::}\NormalTok{by\_name}\OperatorTok{(}\StringTok{"database.org"}\OperatorTok{);}
    \KeywordTok{auto}\NormalTok{ dbdisk }\OperatorTok{=}\NormalTok{ dbhost}\OperatorTok{{-}\textgreater{}}\NormalTok{get\_disks}\OperatorTok{().}\NormalTok{front}\OperatorTok{();}
    \CommentTok{// Get list of compute nodes in my cluster}
    \KeywordTok{auto}\NormalTok{ nodes }\OperatorTok{=}\NormalTok{ this\_actor}\OperatorTok{::}\NormalTok{get\_host}\OperatorTok{()}
    \OperatorTok{{-}\textgreater{}}\NormalTok{get\_englobing\_zone}\OperatorTok{(){-}\textgreater{}}\NormalTok{get\_all\_hosts}\OperatorTok{();}
    \CommentTok{// Create my mailbox}
    \KeywordTok{auto}\NormalTok{ mailbox }\OperatorTok{=}\NormalTok{ Mailbox}\OperatorTok{::}\NormalTok{by\_name}\OperatorTok{(}\NormalTok{this\_actor}\OperatorTok{::}\NormalTok{get\_name}\OperatorTok{());}

    \CommentTok{// Main loop}
    \ControlFlowTok{while} \OperatorTok{(}\KeywordTok{true}\OperatorTok{)} \OperatorTok{\{}
    \CommentTok{// Ask the coordinator for work (32{-}byte message)}
\NormalTok{    Mailbox}\OperatorTok{::}\NormalTok{by\_name}\OperatorTok{(}\StringTok{"coordinator\_mb"}\OperatorTok{){-}\textgreater{}}\NormalTok{put}\OperatorTok{(}\NormalTok{mailbox}\OperatorTok{,} \DecValTok{32}\OperatorTok{);}
    \CommentTok{// Get a workunit back}
    \KeywordTok{auto}\NormalTok{ wu }\OperatorTok{=}\NormalTok{ my\_mailbox}\OperatorTok{{-}\textgreater{}}\NormalTok{get}\OperatorTok{\textless{}}\NormalTok{WorkUnit}\OperatorTok{\textgreater{}();}
    \CommentTok{// If it\textquotesingle{}s a poison pill, terminate}
    \ControlFlowTok{if} \OperatorTok{(}\NormalTok{wu}\OperatorTok{{-}\textgreater{}}\NormalTok{iterations }\OperatorTok{==} \DecValTok{0}\OperatorTok{)} \ControlFlowTok{break}\OperatorTok{;}

    \CommentTok{// Start and MPI program to perform the work}
\NormalTok{    SMPI\_app\_instance\_start}\OperatorTok{(}\NormalTok{this\_actor}\OperatorTok{::}\NormalTok{get\_cname}\OperatorTok{(),}
    \OperatorTok{[}\NormalTok{wu}\OperatorTok{,}\NormalTok{ dbhost}\OperatorTok{,}\NormalTok{ dbdisk}\OperatorTok{])} \OperatorTok{\{}
\NormalTok{    MPI\_Init}\OperatorTok{();}
    \DataTypeTok{int}\NormalTok{ numprocs}\OperatorTok{,}\NormalTok{ myrank}\OperatorTok{;}
\NormalTok{    MPI\_Comm\_size}\OperatorTok{(}\NormalTok{MPI\_COMM\_WORLD}\OperatorTok{,} \OperatorTok{\&}\NormalTok{numprocs}\OperatorTok{);}
\NormalTok{    MPI\_Comm\_rank}\OperatorTok{(}\NormalTok{MPI\_COMM\_WORLD}\OperatorTok{,} \OperatorTok{\&}\NormalTok{myrank}\OperatorTok{);}
    \CommentTok{// Allocate a data buffer}
    \DataTypeTok{void} \OperatorTok{*}\NormalTok{data }\OperatorTok{=}\NormalTok{ SMPI\_SHARED\_MALLOC}\OperatorTok{(}\NormalTok{wu}\OperatorTok{{-}\textgreater{}}\NormalTok{size }\OperatorTok{*}\NormalTok{ numprocs}\OperatorTok{);}
    \ControlFlowTok{for} \OperatorTok{(}\DataTypeTok{int}\NormalTok{ it }\OperatorTok{=} \DecValTok{0}\OperatorTok{;}\NormalTok{ it}\OperatorTok{\textless{}}\NormalTok{ wu}\OperatorTok{{-}\textgreater{}}\NormalTok{iterations}\OperatorTok{;}\NormalTok{ it}\OperatorTok{++)} \OperatorTok{\{}
    \CommentTok{// Perform BSP compute work}
\NormalTok{    this\_actor}\OperatorTok{::}\NormalTok{execute}\OperatorTok{(}\NormalTok{wu}\OperatorTok{{-}\textgreater{}}\NormalTok{work}\OperatorTok{);}
    \CommentTok{// Perform BSP all{-}to{-}all communication}
\NormalTok{    MPI\_Alltoall}\OperatorTok{(}\NormalTok{data}\OperatorTok{,}\NormalTok{ wu}\OperatorTok{{-}\textgreater{}}\NormalTok{size}\OperatorTok{,}\NormalTok{ MPI\_CHAR}\OperatorTok{,}\NormalTok{ data}\OperatorTok{,}\NormalTok{ wu}\OperatorTok{{-}\textgreater{}}\NormalTok{size}\OperatorTok{,}\NormalTok{ MPI\_CHAR}\OperatorTok{,}\NormalTok{ MPI\_COMM\_WORLD}\OperatorTok{);}
    \CommentTok{// Rank 0 performs a database update}
    \ControlFlowTok{if} \OperatorTok{(}\NormalTok{myrank }\OperatorTok{==} \DecValTok{0}\OperatorTok{)} \OperatorTok{\{}
    \CommentTok{// Create 1MB read activity on the database disk}
    \KeywordTok{auto}\NormalTok{ io\_activity }\OperatorTok{=}
\NormalTok{    dbdisk}\OperatorTok{{-}\textgreater{}}\NormalTok{io\_init}\OperatorTok{(}\DecValTok{1}\OperatorTok{*}\NormalTok{MB}\OperatorTok{,}\NormalTok{ Io}\OperatorTok{::}\NormalTok{OpType}\OperatorTok{::}\NormalTok{READ}\OperatorTok{);}
    \CommentTok{// Create 1MB communication activity}
    \CommentTok{// from my host to the database host}
    \KeywordTok{auto}\NormalTok{ comm\_activity }\OperatorTok{=}
\NormalTok{    Comm}\OperatorTok{::}\NormalTok{sendto\_init}\OperatorTok{(}\NormalTok{this\_actor}\OperatorTok{::}\NormalTok{get\_host}\OperatorTok{(),}
\NormalTok{    dbhost}\OperatorTok{)}
    \OperatorTok{{-}\textgreater{}}\NormalTok{set\_payload\_size}\OperatorTok{(}\DecValTok{1}\OperatorTok{*}\NormalTok{MB}\OperatorTok{));}
    \CommentTok{// Create activity dependency}
\NormalTok{    comm\_activity}\OperatorTok{{-}\textgreater{}}\NormalTok{add successor}\OperatorTok{(}\NormalTok{io\_activity}\OperatorTok{);}
    \CommentTok{// Start the I/O activity, but it won\textquotesingle{}t start}
    \CommentTok{// until the communication activity completes io\_activity{-}\textgreater{}start();}
    \CommentTok{// Start the communication activity}
    \CommentTok{// and wait for its completion}
\NormalTok{    comm\_activity}\OperatorTok{{-}\textgreater{}}\NormalTok{start}\OperatorTok{(){-}\textgreater{}}\NormalTok{wait}\OperatorTok{();}
    \OperatorTok{\}}
    \OperatorTok{\}}
    \CommentTok{// Free the data buffer}
\NormalTok{    SMPI\_SHARED\_FREE}\OperatorTok{(}\NormalTok{data}\OperatorTok{);}
\NormalTok{    MPI\_Finalize}\OperatorTok{();}
    \OperatorTok{\},}\NormalTok{ nodes}\OperatorTok{);}

    \CommentTok{// Wait for the MPI program to terminate}
\NormalTok{    SMPI\_app\_instance\_join}\OperatorTok{(}\NormalTok{this\_actor}\OperatorTok{::}\NormalTok{get\_cname}\OperatorTok{());}
    \OperatorTok{\}}
\OperatorTok{\}} 
\end{Highlighting}
\end{Shaded}

Figure 12: Worker C++ functor for the example in Section 5.4.

At line 29, the program allocates sufficient memory for the buffer that
will be used for all-to-all communication, based on the workunit's
specification. While this could be done using a standard malloc(), the
simulation's memory footprint would become too large when simulating the
program's execution with many MPI processes. Instead, the data is
allocated using the SMPI\_SHARED\_MALLOC() macro. This macro, as
explained in details in {[}51{]}, allocates memory that is shared by all
simulated MPI processes. While using this macro would make a real MPI
program incorrect, for a program that does not perform any actual
computation but only simulates the execution of computation volumes, it
makes it possible to simulate the execution of large-scale programs on a
single computer. The MPI program then loops in BSP fashion (lines
30-51). At each iteration of the loop, each process performs some
computation volume as specified by the workunit (line 32), followed by
an MPI all-to-all communication for a data volume also specified by the
workunit (line 34). The process with rank 0 performs extra work at each
iteration (lines 36-50). It creates an I/O activity to be performed on
the database disk (line 38) and a communication activity to be performed
between the host this process runs on and the database host (line 41).
Both these activities have a 1MB payload, and the I/O activity depends
on the communication activity (line 43). The process starts both
activities and waits for the completion of the communication activity
(lines 46 and 49). Finally, once all BSP iterations have been performed,
each simulated MPI process frees the data buffer and calls
MPI\_Finalize() (lines 53-54).

\section{6 Impact on simulation
practice}\label{impact-on-simulation-practice}

In this section we demonstrate the impact of the enhanced usability and
extensibility of SimGrid v4 on simulation practice. The simulation
abstractions and programming models that we have described in the
previous sections allow users to implement simulators for many different
use cases with reasonably low effort. This has led users to develop
simulators for several domains, as seen in Figure 1. Hereafter we
describe notable projects for which SimGrid provides a foundation and
that form a large SimGrid ecosystem. Some of these projects have
produced simulators or simulation frameworks that aim at avoiding
duplication of effort in research communities. Others have produced
simulators or simulation-based tools for various production uses. All
these projects are a testimony to SimGrid's versatility and have all
benefited from its capabilities in terms of accuracy and scalability.
But, in all that follows, we specifically explain how they have
benefited from the extensibility and usability enhancements described in
Sections 4 and 5.

\section{6.1 Distributed cyberinfrastructure
simulation}\label{distributed-cyberinfrastructure-simulation}

Many researchers wish to simulate the execution of various application
workloads on distributed platforms, and often end up re-implementing the
same simulation abstractions and mechanisms. WRENCH {[}53{]} is a
simulation framework that provides implementations of highly
configurable services that users can re-use as building blocks in their
simulators. It also removes the burden of implementing inter-process
communication, which is labor-intensive and error-prone when developing
a simulator of a complex distributed system. The user only writes the
code of one kind of actors called execution controllers. These
controllers interact with the services deployed on the simulated
platform using simple APIs, so as to execute application workloads
defined by data and compute volumes with arbitrary dependencies. As a
result, it is possible to implement simulators of complex deployments
and runtime systems with low software engineering effort. WRENCH
achieves this objective because SimGrid's API is expressive enough to
describe a wide range of interoperable services and sufficiently usable
to render the implementation of these services tractable. In particular,
WRENCH relies on plugins (Section 4.2), on the ability to mix different
programming models (Section 5.4), and on composite activities (Section
5.3).

\section{6.2 Resources and jobs management systems
simulation}\label{resources-and-jobs-management-systems-simulation}

An active research area is Resources and Jobs Management Systems (RJMS),
i.e., the systems in charge of the scheduling of user jobs on shared
parallel computing platforms. In particular, RJMS must employ so-called
batch scheduling algorithms, which have been the subject of active
research for decades, and are typically evaluated in simulation. The
Batsim project {[}54{]} is a SimGrid-based RJMS simulator. While most
research in this area simulates batch jobs at a completely abstract
level (e.g., rectangles in a Gantt chart), Batsim users can describe
workloads in which each job is defined by a profile that encodes
specific computation, communication, and I/O patterns. Job executions
are then simulated so that they generate load, contention, and
electrical power consumption on hardware resources. The ElastiSim
project {[}55{]} shares Batsim's objectives but specifically targets the
simulation of malleable batch jobs that can both adapt their resource
demands and report on their progress at runtime. These two features
dramatically increase the design space for batch scheduling algorithms
as these algorithms can make decisions regarding a job during that job's
execution. Batsim uses the pstate feature for rich resource descriptions
(Section 4.1), and both Batsim and ElastiSim rely on SimGrid's composite
ptask activities (Section 5.3).

\section{6.3 HPC runtimes and
applications}\label{hpc-runtimes-and-applications}

Many researchers have used simulation for research and development in
the field of HPC, a field in which SimGrid has seen a lot of usage (as
seen in the fraction of works that target Cluster platforms in Figure
1). The most common use of simulation if to evaluate the performance of
application workloads when executed at different scales on candidate
platforms. Fast simulations can provide a first performance
approximation before proceeding with resource-intensive testing on or
even purchasing of real compute infrastructures.

Several large projects have used SimGrid in this context due to its SMPI
component being able to simulate the execution of (almost) unmodified
MPI applications. Evaluating the performance of classical application
workloads on particular platforms and/or with particular runtime systems
is so important to the HPC community that several ``proxy apps'' have
been developed and maintained over the years. These correspond to
representative MPI applications for various application domains (ECP
proxy applications {[}56{]}, MeteoFrance proxy applications {[}57{]})
and to standard MPI benchmarks from various benchmark suites (HPL,
CodeVault, Trinity-Nersc, CORAL). Many such applications have been
executed directly with SMPI, and are now part of SimGrid's integration
testing infrastructure {[}58{]}, along with the test suite for the
OpenMPI implementation of the MPI standard {[}59{]} and the Intel MPI
benchmarks {[}60{]}. The S4BXI project {[}61{]} has developed a
full-fledged simulator of the Portals 4 network API {[}62{]} using
SimGrid, in the context of MPI applications. The goal is not only to
enable in-simulation performance evaluation capabilities, but also to
perform what-if analyses to decide promising areas for hardware platform
design optimization. S4BXI uses a multimodel simulation approach, by
which SimGrid's SMPI feature is used for fast simulation of some
portions of the execution and much slower, but highly accurate,
simulation techniques are used to simulate other portions. The BigDFT
project {[}63, 64{]}, which provides open-source software for simulating
macromolecular systems at the nanoscale, uses SimGrid to include
simulated executions in its continuous integration process and
regression testing process. Finally, SimGrid has also been used by Intel
and Bull/Atos to explore different solutions and configurations for
hardware components of HPC platforms, such as network interconnects or
node size {[}3{]}.

Key to the success of the above projects is SimGrid's scalability and
accuracy levels. But key to their feasibility in the first place is the
use of some of the extensibility and usability features added to SimGrid
v4. All these projects use the composite platform features (Section 5.1)
for better usability of SimGrid when describing complex HPC platforms.
The advanced modeling mechanisms (Section 4.3) are also used heavily to
make the accurate simulation models of MPI point-to-point and collective
communications. Some of these projects also rely on the extensibility
plugin feature (Section 4.2), e.g., for energy consumption simulation.

A notable HPC project that uses SimGrid but does not target MPI is
StarPU {[}65{]}. It provides a unified runtime system for programming
heterogeneous multicore architectures (i.e., multicore processors with
accelerators and/or coprocessors, such as GPUs). StarPU integrates with
SimGrid so that StarPU applications can seamlessly execute in simulation
mode, which is a powerful testimony to SimGrid's usability. Simulated
executions are also used for StarPU's continuous integration and
performance evaluation purposes, and have been instrumental in
uncovering performance bugs in StarPU. StarPU builds on SimGrid's models
and abstractions to develop its own models and abstractions for the
simulation of GPUs {[}48{]}, which was made possible by the
extensibility capabilities of SimGrid v4, namely rich and unified
resources (Section 4.1) and advanced modeling (Section 4.3).

\section{7 Impact on scalability}\label{impact-on-scalability}

The usability and extensibility features described in this work have
required a full rewrite of SimGrid v3's monolithic simulation core as
well as performing API overhauls, including moving from C to C++.
Therefore, one may wonder whether these developments have had a negative
impact on scalability. To answer this question we consider a benchmark
simulator that implements a coordinator-worker application (available on
GitHub {[}66{]}). Each worker actor executes on one core of a multi-core
host and a single coordinator actor executes on one core of a separate
host. All hosts are interconnected over a network topology that
comprises 100 network links. The route between any two hosts consists of
10 randomly selected links. The core speed of each host is sampled
uniformly between 100 and 200 Mflop/sec, and the bandwidth of each link
is sampled uniformly between 100 and 200 MB/sec.~The coordinator
greedily assigns workunits to idle workers. Each workunit entails
sending 100 MB of data from the coordinator to the worker, and
performing 100 Mflops of computation at the worker. The simulation of
energy consumption is enabled. We built this benchmark for both SimGrid
v3 and SimGrid v4 in the same Debian 12 Docker image (using the same
compiler), and executed it on one core of a dedicated 2.3GHz Intel Xeon
Platinum 8380 CPU.

Figure 13 shows results when simulating the execution of between 4,000
and 60,000 workunits on 1,000 4-core hosts. As expected, simulation time
increases with the number of workunits (since more discrete events need
to be simulated) and memory footprint is roughly constant (since memory
is freed each time a workunit completes). The most striking observation
is the large reduction in memory footprint when going from SimGrid v3 to
SimGrid v4. Specifically, SimGrid v4 leads to at least a 4.1x reduction
in memory footprint. This is due to optimizations of data structures and
to the move from C to C++, allowing the use of smart pointers with
reference counting

\textit{[Image omitted: images/a2911463ae2c4f2a22851fed67dbe42ceaddbab40183d0b87e77f14f1fe6976e.jpg]}

line

{\def\LTcaptype{none} % do not increment counter
\begin{longtable}[]{@{}lll@{}}
\toprule\noalign{}
\# of workunits & Simulation Time (sec) & Maximum RSS (MB) \\
\midrule\noalign{}
\endhead
\bottomrule\noalign{}
\endlastfoot
0 & 0.5 & 350 \\
10000 & 0.7 & 350 \\
20000 & 1.0 & 350 \\
30000 & 1.3 & 350 \\
40000 & 1.6 & 350 \\
50000 & 2.0 & 350 \\
60000 & 2.3 & 350 \\
\end{longtable}
}

Figure 13: Simulation time and memory footprint vs.~number of workunits
on 1,000 4-core hosts for a benchmark coordinator-worker simulator,
using SimGrid v3 and v4. Lines show average values over 10 repeats;
error bars show minimum and maximum values.

for automated garbage collection to avoid memory leaks. The use of smart
pointers and the extensibility and usability improvements described in
this work do cause increases in simulation time. But these increases are
largely offset by usage of the highly optimized data containers provided
by the C++ standard library and by the rounds of optimizations that have
been applied to the source code over the last decade. As a result, for
this benchmark SimGrid v4 leads to at least a 1.45x reduction in
simulation time when compared to SimGrid v3.

\textit{[Image omitted: images/911424f17eee879ad8affcae215f3d8d5627fc7db525ec898c99cbaabd99af7f.jpg]}

line

{\def\LTcaptype{none} % do not increment counter
\begin{longtable}[]{@{}
  >{\raggedright\arraybackslash}p{(\linewidth - 4\tabcolsep) * \real{0.4478}}
  >{\raggedright\arraybackslash}p{(\linewidth - 4\tabcolsep) * \real{0.3134}}
  >{\raggedright\arraybackslash}p{(\linewidth - 4\tabcolsep) * \real{0.2388}}@{}}
\toprule\noalign{}
\begin{minipage}[b]{\linewidth}\raggedright
\# of cores / host (1000 hosts)
\end{minipage} & \begin{minipage}[b]{\linewidth}\raggedright
Simulation Time (sec)
\end{minipage} & \begin{minipage}[b]{\linewidth}\raggedright
Maximum RSS (MB)
\end{minipage} \\
\midrule\noalign{}
\endhead
\bottomrule\noalign{}
\endlastfoot
0 & 0 & 0 \\
10 & 5 & 1000 \\
20 & 10 & 2000 \\
30 & 15 & 3000 \\
40 & 20 & 4000 \\
50 & 25 & 5000 \\
60 & 30 & 6000 \\
\end{longtable}
}

Figure 14: Simulation time and memory footprint vs.~number of workers
for executing 60,000 workunits for a benchmark coordinator-worker
simulator, using SimGrid v3 and v4. Lines show average values over 10
repeats; error bars show minimum and maximum values.

Figure 14 shows results when simulating the execution of 60,000
workunits on 1,000 hosts where the number of cores per host varies from
4 to 60. As expected memory footprint increases due to the fact that
each worker actor on each core of the simulated worker hosts has a
memory footprint throughout the whole simulated execution. One may
expect the execution time to be constant, since the number of simulated
discrete events does not depend on the number of worker actors. But as
the number of worker actors increases, so does the number of concurrent
simulated activities, which in turn increases the computational
complexity of the LMM solver (linearly). In these results, for the same
reasons as for the results in Figure 13, SimGrid v4 leads to at least a
1.40x memory footprint reduction and a 4.21x simulation time reduction
when compared to SimGrid v3.

Overall, while this work has focused on describing features that make
SimGrid v4 significantly more extensible and more usable than SimGrid
v3, the above results show that it is also significantly more scalable.

\section{8 Conclusion}\label{conclusion}

The primary research and development goals of the SimGrid simulation
framework have been accuracy and scalability. A key accomplishment is
that SimGrid can achieve high accuracy and scalability while remaining
versatile, i.e., it has been used for conducting simulations for broad
ranges of PDC domains and platforms. In spite of these achievements, a
shortcoming of SimGrid has been its usability: the relatively low-level
simulation abstractions it provided made the implementation of
simulators of complex systems a labor-intensive process. Another
shortcoming was extensibility. Although SimGrid provided simulation
models that catered to the common case, it was difficult to users to
extend these models for research- or domain-specific purposes. In this
work we have described several usability and extensibility improvements
implemented over the last decade. These improvements have enabled many
research results and a led to a rich ecosystem of development and
production tools.

Many simulation challenges have been addressed over the last couple of
decades resulting by PDC simulation frameworks that provide various
levels of accuracy, scalability, versatility, extensibility, and
usability. Although future improvements along these five axes are still
possible, the state of PDC simulation is now sufficiently mature for
broader and overarching challenges to be tackled. One such challenge is
simulation calibration: picking the simulation model parameter values in
a way that maximizes simulation accuracy with respect to some
ground-truth. A review of research articles that have used SimGrid-based
simulators in recent years shows that calibration is often not
performed, and that when it is performed is is mostly a labor-intensive
and manual process {[}67{]}. A clear future research direction is the
development of an automated simulation calibration tool that can be used
by PDC simulator users to achieve desirable trade-offs between accuracy
and scalability. Another broad challenge is that of the duplication of
effort when it comes to simulator development. Many SimGrid simulator
developers implement their own simulation abstractions using the
mechanisms in Sections 4 and 5. The question is that of how these
custom-developed abstractions, which would often be useful to others,
can be contributed back to SimGrid. Relying on pull requests is too
labor-intensive (code reviews, increasing amount of

code that must be maintained and documented using the SimGrid standards,
regular release schedule constraints, etc.). An alternative approach,
used successfully by the ns-3 network simulation framework {[}68{]}, is
to establish a SimGrid app store on which contributors can publish their
simulation abstractions as standalone components (composite
abstractions, programmatic platform descriptions, plugin
implementations, implementations of reusable distributed system building
blocks).

\section{References}\label{references}

{[}1{]} The SimGrid Project, http://simgrid.org (2024).\\
{[}2{]} H. Casanova, A. Giersch, A. Legrand, M. Quinson, F. Suter,
Versatile, Scalable, and Accurate Simulation of Distributed Applications
and Platforms, Journal of Parallel and Distributed Computing 74 (10)
(2014) 2899--2917. doi:10.1016/j.jpdc.2014.06.008.\\
{[}3{]} P. Bédaride, A. Degomme, S. Genaud, A. Legrand, G. S.
Markomanolis, M. Quinson, M. Stillwell, F. Suter, B. Videau, Toward
Better Simulation of MPI Applications on Ethernet/TCP Networks, in:
Proceedings of the 4th International Workshop on Performance Modeling,
Benchmarking and Simulation of High Performance Computer Systems (PMBS),
2013. doi:10.1007/978-3-319-10214-6\_8.\\
{[}4{]} P. Velho, L. M. Schnorr, H. Casanova, A. Legrand, On the
Validity of Flow-Level Tcp Network Models for Grid and Cloud
Simulations, ACM Transactions on Modeling and Computer Simulation 23 (4)
(Dec.~2013). doi:10.1145/2517448.\\
{[}5{]} P. Velho, A. Legrand, Accuracy Study and Improvement of Network
Simulation in the SimGrid Framework, in: Proceedings of the 2nd Intl.
Conf. on Simulation Tools and Techniques, 2009.
doi:10.4108/ICST.SIMUTOOLS2009.5592.\\
{[}6{]} K. Fujiwara, H. Casanova, Speed and Accuracy of Network
Simulation in the SimGrid Framework, in: Proceedings of the 1st
International Workshop on Network Simulation Tools, 2007.
doi:10.4108/nstools.2007.2010.\\
{[}7{]} A. Lèbre, A. Legrand, F. Suter, P. Veyre, Adding Storage
Simulation Capacities to the SimGrid Toolkit: Concepts, Models, and API,
in: Proceedings of the 15th IEEE/ACM International Symposium on Cluster,
Cloud and Grid Computing, Shenzen, China, 2015.
doi:10.1109/CCGrid.2015.134.\\
{[}8{]} L. Pouilloux, T. Hirofuchi, A. Lebre, SimGrid VM: Virtual
Machine Support for a Simulation Framework of Distributed Systems, IEEE
Transactions on Cloud Computing (Sep.~2015).
doi:10.1109/TCC.2015.2481422.

{[}9{]} A. Degomme, A. Legrand, G. Markomanolis, M. Quinson, M.
Stillwell, F. Suter, Simulating MPI applications: the SMPI approach,
IEEE Transactions on Parallel and Distributed Systems 18 (8) (2017)
2387--2400. doi:10.1109/TPDS.2017.2669305.\\
{[}10{]} A. Rizvi, T. Toha, M. Lunar, M. Adnan, A. Islam, Cooling Energy
Integration in SimGrid, in: Proceedings of the 2017 International
Conference on Networking, Systems and Security, 2017, pp.~132--137. doi:
10.1109/NSysS.2017.7885814.\\
{[}11{]} F. C. Heinrich, T. Cornebize, A. Degomme, A. Legrand, A.
Carpen-Amarie, S. Hunold, A. Orgerie, M. Quinson, Predicting the
Energy-Consumption of MPI Applications at Scale Using Only a Single
Node, in: Proceedings of the 2017 IEEE International Conference on
Cluster Computing, 2017, pp.~92--102. doi:10.1109/CLUSTER.2017.66.\\
{[}12{]} L. Stanisic, E. Agullo, A. Buttari, A. Guermouche, A. Legrand,
F. Lopez, B. Videau, Fast and Accurate Simulation of Multithreaded
Sparse Linear Algebra Solvers, in: Proceedings of the 2015 IEEE 21st
International Conference on Parallel and Distributed Systems, 2015,
pp.~481--490. doi:10.1109/ICPADS.2015.67.\\
{[}13{]} A. Fanfakh, Predicting the Performance of MPI Applications over
Different Grid Architectures, Journal of University of Babylon for Pure
and Applied Sciences 27 (1) (2019) 468--477.
doi:10.29196/jubpas.v27i1.2232.\\
{[}14{]} L. Stanisic, A Reproducible Research Methodology for Designing
and Conducting Faithful Simulations of Dynamic HPC Applications,
Ph.D.~thesis, Université Grenoble Alpes, France (2015). URL
https://theses.hal.science/tel-01248109v2/document\\
{[}15{]} T. Cornebize, A. Legrand, F. C. Heinrich, Fast and Faithful
Performance Prediction of MPI Applications: the HPL Case Study, in:
Proceedings of the 2019 IEEE International Conference on Cluster
Computing, 2019, pp.~1--11. doi:10.1109/CLUSTER.2019.8891011.\\
{[}16{]} SimGrid's Use in Research Publications,
https://simgrid.org/usages.html (2023).\\
{[}17{]} G. Kecskemeti, S. Ostermann, R. Prodan, Fostering
Energy-Awareness in Simulations Behind Scientific Workflow Management
Systems, in: Proceedings of the 7th IEEE/ACM Intl. Conf. on Utility and
Cloud Computing, 2014, pp.~29--38. doi:10.1109/UCC.2014.11.\\
{[}18{]} G. F. Riley, T. R. Henderson, The ns-3 Network Simulator,
Springer Berlin Heidelberg, Berlin, Heidelberg, 2010, pp.~15--34.
doi:10.1007/978-3-642-12331-3\_2.

{[}19{]} L. Mészáros, A. Varga, M. Kirsche, INET Framework, Springer
International Publishing, Cham, 2019, pp.~55--106.
doi:10.1007/978-3-030-12842-5\_2.\\
{[}20{]} N. Binkert, B. Beckmann, G. Black, S. K. Reinhardt, A. Saidi,
A. Basu, J. Hestness, D. R. Hower, T. Krishna, S. Sardashti, R. Sen, K.
Sewell, M. Shoaib, N. Vaish, M. D. Hill, D. A. Wood, The gem5 simulator,
SIGARCH Comput. Archit. News 39 (2) (2011) 1--7.
doi:10.1145/2024716.2024718.\\
{[}21{]} T. E. Carlson, W. Heirman, L. Eeckhout, Sniper: exploring the
level of abstraction for scalable and accurate parallel multi-core
simulation, in: Proc. of the 2011 International Conference for High
Performance Computing, Networking, Storage and Analysis, Association for
Computing Machinery, New York, NY, USA, 2011.
doi:10.1145/2063384.2063454.\\
{[}22{]} D. Gouk, M. Kwon, J. Zhang, S. Koh, W. Choi, N. S. Kim, M.
Kandemir, M. Jung, Amber: Enabling Precise Full-System Simulation with
Detailed Modeling of All SSD Resources, in: 2018 51st Annual IEEE/ACM
International Symposium on Microarchitecture (MICRO), 2018,
pp.~469--481. doi:10.1109/MICRO.2018.00045.\\
{[}23{]} G. G. Castañé, A. Núñez, J. Carretero, iCanCloud: A Brief
Architecture Overview, in: 2012 IEEE 10th International Symposium on
Parallel and Distributed Processing with Applications, 2012,
pp.~853--854. doi: 10.1109/ISPA.2012.131.\\
{[}24{]} D. Kliazovich, P. Bouvry, Y. Audzevich, S. U. Khan, GreenCloud:
A Packet-Level Simulator of Energy-Aware Cloud Computing Data Centers,
in: 2010 IEEE Global Telecommunications Conference GLOBECOM 2010, 2010,
pp.~1--5. doi:10.1109/GLOCOM.2010.5683561.\\
{[}25{]} R. M. Fujimoto, Parallel discrete event simulation, Commun. ACM
33 (10) (1990) 30--53. doi:10.1145/84537.84545.\\
{[}26{]} J. Cope, N. Liu, S. Lang, P. Carns, C. Carothers, R. Ross,
CODES: Enabling Co-Design of Multilayer Exascale Storage Architectures,
in: Proc. of the Workshop on Emerging Supercomputing Technologies,
2011.\\
{[}27{]} C. Carothers, D. Bauer, S. Pearce, ROSS: A High-Performance,
Low Memory, Modular Time Warp System, in: Proc. of the 14th ACM/IEEE/SCS
Workshop of Parallel on Distributed Simulation, 2000, pp.~53--60.
doi:10.1109/PADS.2000.847144.\\
{[}28{]} S. Böhm, C. Engelmann, xSim: The extreme-scale simulator, in:
Proc. of the International Conference on High Performance Computing \&
Simulation, 2011, pp.~280--286. doi:10.1109/HPCSim.2011.5999835.

{[}29{]} M.-Y. Hsieh, R. Riesen, K. Thompson, W. Song, A. Rodrigues,
SST: A Scalable Parallel Framework for Architecture-Level Performance,
Power, Area and Thermal Simulation, The Computer Journal 55 (2) (2012)
181--191. doi:10.1093/comjnl/bxr069.\\
{[}30{]} SST/macro 14.1: User's Manual,
https://raw.githubusercontent.com/sstsimulator/sst-macro/refs/heads/master/manual-sstmacro-14.1.pdf
(2024).\\
{[}31{]} R. Buyya, M. Murshed, GridSim: A Toolkit for the Modeling and
Simulation of Distributed Resource Management and Scheduling for Grid
Computing, Concurrency and Computation: Practice and Experience 14 (11
2002). doi:10.1002/cpe.710.\\
{[}32{]} T. Goyal, A. Singh, A. Agrawal, Cloudsim: Simulator for Cloud
Computing Infrastructure and Modeling, Procedia Engineering 38 (2012)
3566--3572. doi:https://doi.org/10.1016/j.proeng.2012.06.412.\\
{[}33{]} E. U. Yousuf Khan, T. Rahim Soomro, M. Nawaz Brohi, iFogSim: A
Tool for Simulating Cloud and Fog Applications, in: Proceedings of the
International Conference on Cyber Resilience, 2022, pp.~01--05. doi:
10.1109/ICCR56254.2022.9996018.\\
{[}34{]} G. Kecskemeti, DISSECT-CF: A Simulator to Foster Energy-Aware
Scheduling in Infrastructure Clouds, Simulation Modelling Practice and
Theory 58 (2015) 188--218.
doi:https://doi.org/10.1016/j.simpat.2015.05.009.\\
{[}35{]} S. Ostermann, K. Plankensteiner, R. Prodan, T. Fahringer,
GroudSim: An Event-Based Simulation Framework for Computational Grids
and Clouds, in: Proceedings of the Euro-Par 2010 Parallel Processing
Workshops, 2011, pp.~305--313. doi:10.1007/978-3-642-21878-1\_38.\\
{[}36{]} F. Mastenbroek, G. Andreadis, S. Jounaid, W. Lai, J. Burley, J.
Bosch, E. van Eyk, L. Versluis, V. van Beek, A. Iosup, Opendc 2.0:
Convenient modeling and simulation of emerging technologies in cloud
datacenters, in: Proc. of the 21st IEEE/ACM International Symposium on
Cluster, Cloud and Internet Computing (CCGrid), 2021, pp.~455--464.
doi:10.1109/CCGrid51090.2021.00055.\\
{[}37{]} G. Keller, M. Tighe, H. Lutfiyya, M. Bauer, DC-Sim: A data
centre simulation tool, in: Proc. of the IFIP/IEEE International
Symposium on Integrated Network Management, 2013, pp.~1090--1091.\\
{[}38{]} X. Li, X. Jiang, P. Huang, K. Ye, DartCSim: An enhanced
user-friendly cloud simulation system based on CloudSim with better
performance, in: Proc. of the

2nd IEEE International Conference on Cloud Computing and Intelligence
Systems, Vol. 01, 2012, pp.~392--396. doi:10.1109/CCIS.2012.6664434.\\
{[}39{]} S. Sotiriadis, N. Bessis, N. Antonopoulos, A. Anjum, SimIC:
Designing a New Inter-cloud Simulation Platform for Integrating
Large-Scale Resource Management, in: Proc. of the 27th IEEE
International Conference on Advanced Information Networking and
Applications (AINA), 2013, pp.~90--97. doi:10.1109/AINA.2013.123.\\
{[}40{]} S. Sotiriadis, N. Bessis, N. Antonopoulos, Towards Inter-cloud
Simulation Performance Analysis: Exploring Service-Oriented Benchmarks
of Clouds in SimIC, in: Proc. of the 27th International Conference on
Advanced Information Networking and Applications Workshops, 2013,
pp.~765--771. doi:10.1109/WAINA.2013.196.\\
{[}41{]} Y. Shi, X. Jiang, K. Ye, An Energy-Efficient Scheme for Cloud
Resource Provisioning Based on CloudSim, in: Proc. of the IEEE
International Conference on Cluster Computing, 2011, pp.~595--599.
doi:10.1109/CLUSTER.2011.63.\\
{[}42{]} L. Bobelin, A. Legrand, D. A. G. Márquez, P. Navarro, M.
Quinson, F. Suter, C. Thiery, Scalable Multi-Purpose Network
Representation for Large Scale Distributed System Simulation, in:
Proceedings of the 12th IEEE/ACM International Symposium on Cluster,
Cloud and Grid Computing, 2012, pp.~220--227.
doi:10.1109/CCGrid.2012.31.\\
{[}43{]} M. Quinson, C. Rosa, C. Thiery, Parallel Simulation of
Peer-to-Peer Systems, in: Proceedings of the 2012 12th IEEE/ACM
International Symposium on Cluster, Cloud and Grid Computing, 2012,
pp.~668---675.\\
{[}44{]} R. Jhala, R. Majumdar, Software model checking, ACM Comput.
Surv. 41 (4) (2009). doi:10.1145/1592434.1592438.\\
{[}45{]} M. Laurent, E. Saillard, M. Quinson, The MPI Bugs Initiative: a
Framework for MPI Verification Tools Evaluation, in: Proc. of the 5th
IEEE/ACM International Workshop on Software Correctness for HPC
Applications (Correctness), 2021, pp.~1--9.
doi:10.1109/Correctness54621.2021.00008.\\
{[}46{]} G. Cooperman, M. Quinson, Sthread: In-Vivo Model Checking of
Multithreaded Programs, The Art, Science, and Engineering of Programming
4 (3) (2020). doi:10.22152/programming-journal.org/2020/4/13.\\
{[}47{]} B. Camus, A.-C. Orgerie, M. Quinson, Co-simulation of FMUs and
Distributed Applications with SimGrid, in: Proceedings of the ACM SIGSIM
Conference on Principles of Advanced Discrete Simulation, 2018,
pp.~145--156. doi:10.1145/3200921.3200932.

{[}48{]} L. Stanisic, S. Thibault, A. Legrand, B. Videau, J.-F. Méhaut,
Faithful Performance Prediction of a Dynamic Task-Based Runtime System
for Heterogeneous Multi-Core Architectures, Concurrency and Computation:
Practice and Experience 27 (16) (2015) 4075--4090.
doi:https://doi.org/10.1002/cpe.3555.\\
{[}49{]} C. Courageux-Sudan, L. Guegan, A.-C. Orgerie, M. Quinson, A
Flow-Level Wi-Fi Model for Large Scale Network Simulation, in:
Proceedings of the International Conference on Modeling, Analysis and
Simulation of Wireless and Mobile Systems, 2022.
doi:10.1145/3551659.3559022.\\
{[}50{]} Message Passing Interface Forum, MPI: A message-passing
interface standard version 4.0,
https://www.mpi-forum.org/docs/mpi-4.0/mpi40-report.pdf (Jun.~2021).\\
{[}51{]} P.-N. Clauss, M. Stillwell, S. Genaud, F. Suter, H. Casanova,
M. Quinson, Single Node On-Line Simulation of MPI Applications with
SMPI, in: Proceedings of the 25th IEEE International Parallel and
Distributed Processing Symposium, 2011. doi:10.1109/IPDPS.2011.69.\\
{[}52{]} SimGrid ``Frankenstein'' example simulator,
https://github.com/henricasanova/simgrid\_frankenstein (2024).\\
{[}53{]} H. Casanova, R. Ferreira da Silva, R. Tanaka, S. Pandey, G.
Jethwani, S. Albrecht, J. Oeth, F. Suter, Developing Accurate and
Scalable Simulators of Production Workflow Management Systems with
WRENCH, Future Generation Computer Systems 112 (2020) 162--175.
doi:10.1016/j.future.2020.05.030.\\
{[}54{]} P.-F. Dutot, M. Mercier, M. Poquet, O. Richard, Bat-sim: a
Realistic Language-Independent Resources and Jobs Management Systems
Simulator, in: Proceedings of the 20th Workshop on Job Scheduling
Strategies for Parallel Processing, 2016.
doi:10.1007/978-3-319-61756-5\_10.\\
{[}55{]} T. Özden, T. Beringer, A. Mazaheri, H. M. Fard, F. Wolf,
ElastiSim: A Batch-System Simulator for Malleable Workloads, in:
Proceedings of the 51st International Conference on Parallel Processing,
2023. doi: 10.1145/3545008.3545046.\\
{[}56{]} ECP Proxy Apps,
https://proxyapps.exascaleproject.org/ecp-proxy-apps-suite/ (2023).\\
{[}57{]} Y. Zheng, P. Marguinaud, Simulation of the Performance and
Scalability of Message Passing Interface (MPI) Communications of
Atmospheric Models Running on Exascale Supercomputers, Geoscientific
Model Development 11 (8) (2018) 3409--3426.
doi:10.5194/gmd-11-3409-2018.

{[}58{]} SMPI Proxy Apps, https://framagit.org/simgrid/SMPI-proxy-apps
(2023).\\
{[}59{]} OpenMPI Test Suite, https://github.com/open-mpi/mpi-test-suite
(2023).\\
{[}60{]} Intel®MPI Benchmarks, https://github.com/intel/mpi-benchmarks
(2023).\\
{[}61{]} J. Emmanuel, M. Moy, L. Henrio, G. Pichon, S4BXI: the MPI-ready
Portals 4 Simulator, in: Proceedings of the 29th IEEE International
Symposium on the Modeling, Analysis, and Simulation of Computer and
Telecommunication Systems, 2021, pp.~1--8.
doi:10.1109/MASCOTS53633.2021.9614285.\\
{[}62{]} B. Barrett, R. B. Brightwell, R. Grant, K. Pedretti, K.
Wheeler, K. D. Underwood, R. Riesen, A. B. Maccabe, T. Hudson, S.
Hemmert, The Portals 4.1 Network Programming Interface, Tech.
Rep.~SAND2017-3825, Sandia National Laboratory, Albuquerque, NM
(Apr.~2017). doi:10.2172/1365498.\\
{[}63{]} L. E. Ratcliff, W. Dawson, G. Fisicaro, D. Caliste, S. Mohr, A.
Degomme, B. Videau, V. Cristiglio, M. Stella, M. D'Alessandro, S.
Goedecker, T. Nakajima, T. Deutsch, L. Genovese, Flexibilities of
Wavelets as a Computational Basis Set for Large-Scale Electronic
Structure Calculations, The Journal of Chemical Physics 152 (19) (2020)
194110. doi:10.1063/5.0004792.\\
{[}64{]} F. Affinito, U. Alekseeva, C. Cavazzoni, A. Degomme, P. D.
Delugas, A. Ferretti, A. Garcia, A. Kozhevnikov, P. Ordejón, N.
Spallanzani, Second Report on Code Profiling and Bottleneck
Identification, Deliverable d4.3, European Centre of Excellence in
materials modelling, simulations and design (2018). URL
https://www.max-centre.eu/sites/default/files/D4.3\%20Second\%20report\%20on\%20code\%20profiling\%20and\%20bottleneck\%20identification.pdf\\
{[}65{]} C. Augonnet, S. Thibault, R. Namyst, P.-A. Wacrenier, StarPU: A
Unified Platform for Task Scheduling on Heterogeneous Multicore
Architectures, CCPE - Concurrency and Computation: Practice and
Experience, Special Issue: Euro-Par 2009 23 (2011) 187--198.
doi:10.1002/cpe.1631.\\
{[}66{]} SimGrid ``Coordinator-Worker flashback'' benchmark simulator,
https://github.com/simgrid/coordinator\_worker\_flashback (2024).\\
{[}67{]} J. McDonald, M. Horzela, F. Suter, H. Casanova, Automated
Calibration of Parallel and Distributed Computing Simulators: A Case
Study, in: Proc. of the 25th IEEE International Workshop on Parallel and
Distributed Scientific and Engineering Computing (PDSEC), 2024.

{[}68{]} The ns-3 App Store,
https://www.nsnam.org/docs/contributing/html/external.html (2024).

\end{document}
